\documentclass[sigconf, nonacm, balance]{acmart}
\settopmatter{printacmref=false}

% get rid of warning
\setlength{\headheight}{15.45749pt}

% turn the "country missing" check into a warning instead of an error
\makeatletter
\def\@ACM@checkaffil{%
  \if@ACM@instpresent\else
    \ClassWarningNoLine{\@classname}{No institution present for an affiliation}%
  \fi
  \if@ACM@citypresent\else
    \ClassWarningNoLine{\@classname}{No city present for an affiliation}%
  \fi
  \if@ACM@countrypresent\else
    \ClassWarningNoLine{\@classname}{No country present for an affiliation}%
  \fi
}
\makeatother

% language and encoding
\usepackage[english]{babel}
\usepackage[utf8]{inputenc}
\usepackage[T1]{fontenc} % set font encoding so that special characters are displayed correctly

\usepackage{amsmath, amsfonts} % amssymb
\usepackage{balance}
\usepackage{algorithmic}
\usepackage{algorithm}
\usepackage{array}
\usepackage[caption=false,font=normalsize,labelfont=sf,textfont=sf]{subfig}
\usepackage{textcomp}
\usepackage{stfloats}
\usepackage{url}
\usepackage{verbatim}
\usepackage{graphicx}
%\usepackage{cite}
\usepackage{natbib}

\hyphenation{op-tical net-works semi-conduc-tor IEEE-Xplore}

\setlength{\marginparwidth}{2cm}
\usepackage{hyperref}
\usepackage{csquotes}
\usepackage{tikz}
\usetikzlibrary{shapes.geometric, arrows, fit}

\usepackage{colortbl}
\usepackage{xcolor}
\usepackage{multirow}
\usepackage{booktabs}
\usepackage{xspace}
\usepackage{dashrule}
\usepackage[super]{nth}

\renewcommand{\sectionautorefname}{Section}
\renewcommand{\subsectionautorefname}{Section}
\renewcommand{\subsubsectionautorefname}{Section}

% editorial commands
\newcommand{\todo}[1]{{\textbf{TODO:}\ \textit{#1}}} % command for TODOs
%\newcommand{\comment}[1]{{\textbf{Comment:}\ \textit{#1}}} % command for review comments

% RFC 2119 (https://www.rfc-editor.org/rfc/rfc2119)
% MUST: absolute requirement
\newcommand{\must}{\textsc{must}\xspace}
% MUST NOT: absolute prohibition
\newcommand{\mustnot}{\textsc{must not}\xspace}
% SHOULD: there may exist valid reasons in particular circumstances to ignore a  particular item, but the full implications must be understood and carefully weighed before choosing a different course
\newcommand{\should}{\textsc{should}\xspace}
% SHOULD NOT: there may exist valid reasons in particular circumstances when the particular behavior is acceptable or even useful, but the full implications should be understood and the case carefully weighed before implementing any behavior described with this label
\newcommand{\shouldnot}{\textsc{should not}\xspace}
% MAY: an item is truly optional
\newcommand{\may}{\textsc{may}\xspace}

% command to indicate where certain information should be reported
\newcommand{\paper}{\emph{paper}\xspace}
\newcommand{\supplementarymaterial}{\emph{supplementary material}\xspace}

% configure enumerate/itemize
\usepackage[inline]{enumitem}

% Get author names despite the fact that IEEEtran doesn't support it (it can't be used together with natbib)
%\usepackage[style=ieee, backend=biber,sorting=nty, mincitenames=1, maxcitenames=2]{biblatex}
%\addbibresource{literature.bib}

\newcommand*{\enq}[1]{\enquote{{\itshape#1}}}

% prevent single lines at the start of a paragraph (Schusterjungen)
\clubpenalty = 10000
% prevent single lines at the end of a paragraph (Hurenkinder)
\widowpenalty = 10000 \displaywidowpenalty = 10000

% commands to reference sections

% scope
\newcommand{\scope}{\hyperref[sec:scope]{\emph{Scope}}\xspace}

% study types
\newcommand{\studytypes}{\hyperref[sec:study-types]{\emph{Study Types}}\xspace}
\newcommand{\llmsforresearcher}{\hyperref[sec:llms-as-tools-for-software-engineering-researchers]{\emph{LLMs as Tools for Software Engineering Researchers}}\xspace}
\newcommand{\annotators}{\hyperref[sec:llms-as-annotators]{\emph{LLMs as Annotators}}\xspace}
\newcommand{\judges}{\hyperref[sec:llms-as-judges]{\emph{LLMs as Judges}}\xspace}
\newcommand{\synthesis}{\hyperref[sec:llms-for-synthesis]{\emph{LLMs for Synthesis}}\xspace}
\newcommand{\subjects}{\hyperref[sec:llms-as-subjects]{\emph{LLMs as Subjects}}\xspace}
\newcommand{\llmsforengineers}{\hyperref[sec:llms-as-tools-for-software-engineers]{\emph{LLMs as Tools for Software Engineers}}\xspace}
\newcommand{\llmusage}{\hyperref[sec:studying-llm-usage-in-software-engineering]{\emph{Studying LLM Usage in Software Engineering}}\xspace}
\newcommand{\newtools}{\hyperref[sec:llms-for-new-software-engineering-tools]{\emph{LLMs for New Software Engineering Tools}}\xspace}
\newcommand{\benchmarkingtasks}{\hyperref[sec:benchmarking-llms-for-software-engineering-tasks]{\emph{Benchmarking LLMs for Software Engineering Tasks}}\xspace}

%guidelines
\newcommand{\guidelines}{\hyperref[sec:guidelines]{\emph{Guidelines}}\xspace}
\newcommand{\usagerole}{\hyperref[sec:declare-llm-usage-and-role]{\emph{Declare LLM Usage and Role}}\xspace}
\newcommand{\modelversion}{\hyperref[sec:report-model-version-configuration-and-customizations]{\emph{Report Model Version, Configuration, and Customizations}}\xspace}
\newcommand{\toolarchitecture}{\hyperref[sec:report-tool-architecture-beyond-models]{\emph{Report Tool Architecture beyond Models}}\xspace}
\newcommand{\humanvalidation}{\hyperref[sec:use-human-validation-for-llm-outputs]{\emph{Use Human Validation for LLM Outputs}}\xspace}
\newcommand{\prompts}{\hyperref[sec:report-prompts-their-development-and-interaction-logs]{\emph{Report Prompts, their Development, and Interaction Logs}}\xspace}
\newcommand{\openllm}{\hyperref[sec:use-an-open-llm-as-a-baseline]{\emph{Use an Open LLM as a Baseline}}\xspace}
\newcommand{\benchmarksmetrics}{\hyperref[sec:use-suitable-baselines-benchmarks-and-metrics]{\emph{Use Suitable Baselines, Benchmarks, and Metrics}}\xspace}
\newcommand{\limitationsmitigations}{\hyperref[sec:report-limitations-and-mitigations]{\emph{Report Limitations and Mitigations}}\xspace}

% custom format for subsubsections
\newcommand{\guidelinesubsubsection}[1]{\subsubsection{\bfseries #1}}

\usepackage{listings}
\usepackage{textcomp} % fix single quotes in listings environments
\lstset{
  basicstyle=\ttfamily\scriptsize,
  showspaces=false,
  showstringspaces=false,
  showtabs=false,
  tabsize=2,
  breaklines=true,
  breakatwhitespace=true,
  aboveskip=0.5\baselineskip,
  belowskip=0.5\baselineskip,
  xleftmargin=0.5\parindent,
  xrightmargin=0pt,
  keywordstyle=\color{blue},
  identifierstyle=\color{black},
  commentstyle=\color{gray},
  stringstyle=\color{gray},
  emphstyle=\color{red},
  numbers=left,
  numbersep=8pt,
  upquote=true
%  columns=flexible,
%  numberstyle=\tiny,
%  keepspaces=true
}

\lstnewenvironment{plain} {\lstset{
  numbers=none
}}{}

% boxes for results
\usetikzlibrary{shadows}
\definecolor{framebg}{gray}{0.95}
\usepackage[framemethod=tikz]{mdframed}
\newmdenv[
  leftmargin=0pt,
  rightmargin=0pt,
  usetwoside=false,
  innerleftmargin=6pt,
  innerrightmargin=6pt,
  innertopmargin=4pt,
  innerbottommargin=4pt,
  skipabove=0.5\baselineskip plus 1pt,
  skipbelow=0.5\baselineskip plus 2pt,
  linecolor=black,
  backgroundcolor=framebg,
  linewidth=2pt,
  roundcorner=0pt,
  nobreak=false,
  shadow=false,
  topline=false,
  bottomline=false,
  rightline=false,
  leftline=true,
  font=\itshape
]{framed}

\title{Evaluation Guidelines for Empirical Studies in Software Engineering involving LLMs}

\begin{document}

\author{Sebastian Baltes}
\affiliation{\institution{University of Bayreuth, Germany}}
\email{sebastian.baltes@uni-bayreuth.de}

\author{Florian Angermeir}
\affiliation{\institution{fortiss, Germany}}
\affiliation{\institution{BTH, Sweden}}
\email{angermeir@fortiss.org}

\author{Chetan Arora}
\affiliation{\institution{Monash University, Australia}}
\email{chetan.arora@monash.edu}

\author{Marvin Muñoz Barón\\ Chunyang Chen}
\affiliation{\institution{TU Munich, Germany}}
\email{{marvin.munoz-baron, chun-yang.chen}@tum.de}

\author{Lukas Böhme}
\affiliation{\institution{Hasso-Plattner-Institut, Germany}}
\affiliation{\institution{University of Potsdam, Germany}}
\email{lukas.boehme@hpi.de}

\author{Fabio Calefato}
\affiliation{\institution{University of Bari, Italy}}
\email{fabio.calefato@uniba.it}

\author{Neil Ernst}
\affiliation{\institution{University of Victoria, Canada}}
\email{nernst@uvic.ca}

\author{Davide Falessi}
\affiliation{\institution{University of Rome Tor Vergata, Italy}}
\email{falessi@ing.uniroma2.it}

\author{Brian Fitzgerald}
\affiliation{\institution{Lero, Ireland}}
\affiliation{\institution{University of Limerick, Ireland}}
\email{brian.fitzgerald@ul.ie}

\author{Davide Fucci}
\affiliation{\institution{BTH, Sweden}}
\email{davide.fucci@bth.se}

\author{Marcos Kalinowski}
\affiliation{\institution{PUC Rio de Janeiro, Brazil}}
\email{kalinowski@inf.puc-rio.br}

\author{Stefano Lambiase\\Daniel Russo}
\affiliation{\institution{Aalborg University, Denmark}}
\email{{stla, daniel.russo}@cs.aau.dk}

\author{Mircea Lungu}
\affiliation{\institution{IT University Copenhagen, Denmark}}
\email{mlun@itu.dk}

\author{Lutz Prechelt}
\affiliation{\institution{Freie Universität Berlin, Germany}}
\email{prechelt@inf.fu-berlin.de}

\author{Paul Ralph}
\affiliation{\institution{Dalhousie University, Canada}}
\email{paulralph@dal.ca}

\author{Christoph Treude}
\affiliation{\institution{SMU, Singapore}}
\email{ctreude@smu.edu.sg}

\author{Stefan Wagner}
\affiliation{\institution{TU Munich, Germany}}
\email{stefan.wagner@tum.de}

\renewcommand{\shortauthors}{Baltes et al.}

\begin{abstract}
Large language models (LLMs) are increasingly being integrated into software engineering (SE) research and practice, yet their non-determinism, opaque training data, and evolving architectures complicate the reproduction and replication of empirical studies.
We present a community effort to scope this space, introducing a taxonomy of LLM-based study types together with eight guidelines for designing and reporting empirical studies involving LLMs.
The guidelines present essential (\must) criteria as well as desired (\should) criteria and target transparency throughout the research process.
Our recommendations, contextualized by our study types, are: (1) to declare LLM usage and role; (2) to report model versions, configurations, and fine-tuning; (3) to document tool architectures; (4) to disclose prompts and interaction logs; (5) to use human validation; (6) to employ an open LLM as a baseline; (7) to use suitable baselines, benchmarks, and metrics; and (8) to openly articulate limitations and mitigations.
Our goal is to enable reproducibility and replicability despite LLM-specific barriers to open science.
We maintain the study types and guidelines online as a living resource for the community to use and shape (\href{https://llm-guidelines.org/}{llm-guidelines.org}).
\end{abstract}

%\keywords{software engineering, large language models, empirical research, meta-science, guidelines}

\maketitle

\section{Introduction}
\label{sec:introduction}

In the short period since the release of ChatGPT in November 2022, large language models (LLMs) have changed the software engineering (SE) research landscape.
Although there are numerous opportunities to use LLMs to support SE research and development tasks, solid science needs rigorous empirical evaluations to explore the effectiveness, performance, and robustness of using LLMs to automate different research and development tasks.
LLMs can, for example, be used to support literature reviews, reduce development time, and generate software documentation.
However, it is often unclear how valid, reproducible, and replicable empirical studies involving LLM are.
This uncertainty poses significant challenges for researchers and practitioners who seek to draw reliable conclusions from empirical studies.

The importance of open science practices and documenting the study setup is not to be underestimated~\cite{DBLP:journals/corr/abs-2412-17859}.
One of the primary risks in creating irreproducible and irreplicable results based on studies involving LLMs stems from the variability in model performance due to their inherent non-determinism, but also due to differences in configuration, training data, model architecture, or evaluation metrics.
Slight changes can lead to significantly different results.
The lack of standardized benchmarks and evaluation protocols further hinders the reproducibility and replicability of study results.
Without detailed information on the exact study setup, benchmarks, and metrics used, it is challenging for other researchers to replicate results using different models and tools.
These issues highlight the need for clear guidelines and best practices for designing and reporting studies involving LLMs.
While the SE research community has developed guidelines for conducting and reporting specific types of empirical studies such as controlled experiments (e.g., \emph{Experimentation in Software Engineering}~\cite{DBLP:books/sp/WohlinRHORW24}, \emph{Guide to Advanced Empirical Software Engineering}~\cite{DBLP:books/sp/08/SSS2008}) or their replications (e.g., \emph{A Procedure and Guidelines for Analyzing Groups of Software Engineering Replications}~\cite{DBLP:journals/tse/SantosVOJ21}),  we believe that LLMs have specific intrinsic characteristics that require specific guidelines for researchers to achieve an acceptable level of reproducibility and replicability (see also our previous position paper~\cite{DBLP:conf/wsese/0001BFB25}).
For example, even if we knew the specific version of a commercial LLM used for an empirical study, the reported task performance could still change over time, since commercial models are known to evolve beyond version identifiers~\cite{DBLP:journals/corr/abs-2307-09009}.
Moreover, commercial providers do not guarantee the availability of old model or tool versions indefinitely.
In addition to version differences, LLM performance varies widely depending on configured parameters such as temperature.
Therefore, not reporting the parameter settings severely impacts reproducibility.
Even for ``open'' models such as \emph{Llama}, we do not know how they were fine-tuned for specific tasks and what the exact training data was~\cite{Gibney2024}.
A general problem when evaluating LLMs' performance is that we do not know whether the solution to a certain problem was part of the training data or not.

So far, there are no holistic guidelines for conducting and reporting studies involving LLMs in SE research.
With this community effort, we try to fill this gap.
After outlining our \scope, we continue by introducing a taxonomy of \studytypes before presenting eight \guidelines that the authors of this article co-developed.
The most recent version is always available online.\footnote{\url{https://llm-guidelines.org}}
Other researchers can suggest changes via a public GitHub repository.\footnote{\url{https://github.com/se-uhd/llm-guidelines-website}}

\section{Scope}
\label{sec:scope}

First, we want to clarify that our focus is on LLMs, that is, natural language use cases.
Multi-modal foundational models are beyond the scope of our study types and guidelines.
We are aware that these foundational models have great potential to support software engineering research and practice.
However, due to the diversity of artifacts that can be generated or used as input (e.g., images, audio, and video) and the more demanding hardware requirements, we deliberately focus on LLMs only.
However, our guidelines could be extended in the future to include foundational models beyond natural language text.

Second, given the exponential growth in LLM usage across all research domains, we also want to define the research contexts in which our guidelines apply.
LLMs are already widely used to support several aspects of the overall research process, from fairly simple tasks such as proof-reading, spell-checking, and text translation, to more complex activities such as data coding and synthesis of literature reviews.
The \studytypes and \guidelines we describe are tailored to software engineering (SE) research, but we expect many of our study types to generalize beyond that domain.
On the practical side, we focus on AI for software engineering (AI4SE), that is, studying the support and automation of SE tasks with the help of artificial intelligence (AI), more specifically LLMs (see Section~\llmsforengineers).
In terms of research support, we focus on empirical SE research supported by LLMs (see Section~\llmsforresearcher).
By research support, we mean the active involvement of LLMs in data collection, processing, or analysis.
We consider LLMs supporting the study design or the writing process to be out of scope.

Third, our guidelines mainly target researchers planning, designing, or conducting empirical studies involving LLMs.
Although researchers who review scientific articles written by others can also use our guidelines, for example, to check whether the authors adhere to the essential \must requirements, reviewers are not our main target audience.

\section{Methodology}
\label{sec:methodology}

The starting point for the development of these guidelines was the 2024 meeting of the \emph{International Software Engineering Research Network} (ISERN) in Barcelona, Spain.
The first and last author of this paper organized a session to discuss the need to develop guidelines for empirical studies involving LLMs.
These discussions resulted in a position paper published at the \emph{\nth{2} International Workshop on Methodological Issues with Empirical Studies in Software Engineering} (WSESE 2025)~\cite{DBLP:conf/wsese/0001BFB25}.
The first and last author presented a preprint of that paper at the \emph{\nth{2} Copenhagen Symposium on Human-Centered Software Engineering AI} (CHASEAI 2024), forming a workstream to lead the collaborative refinement and extension of the preliminary guidelines presented in the position paper.

We held bi-weekly meetings, in which we decided on the structure of the study type taxonomy and the guidelines.
During these discussions, we added a new study type (\synthesis), refined the guidelines to \modelversion and to \toolarchitecture, extended the guideline to \prompts, and added the new guidelines to \benchmarksmetrics and to \limitationsmitigations.
Then, at least two co-authors were assigned to each study type and guideline.
They refined and extended the content to align with the intended structure, followed by a peer review and refinement phase.
Afterwards, the first author reviewed all sections and discussed potential changes with the topic leads.
After the internal review was completed, we invited external experts in empirical software engineering to review the study types and guidelines.
After incorporating their feedback,\footnote{The discussion is documented in this GitHub issue: \url{https://github.com/se-uhd/llm-guidelines-website/issues/89}} we compiled this article and published it online.

\section{Study Types}
\label{sec:study-types}

The development of empirical guidelines for software engineering (SE) studies involving LLMs is crucial to ensuring the validity and reproducibility of results.
However, such guidelines must be tailored to different study types that each pose unique challenges.
%Therefore, understanding the classification of these studies is essential for developing appropriate guidelines.
Therefore, we developed a taxonomy of study types that we then use to contextualize the recommendations we provide.
Note that our \guidelines refer to the \studytypes but not the other way around.

Each study type section starts with a \textbf{description}, followed by \textbf{examples} from the SE research community and beyond as well as the \textbf{advantages} and \textbf{challenges} of using LLMs for the respective study type.
The remainder of this section is structured as follows:

\par\noindent\rule{\columnwidth}{0.5pt}
\begin{enumerate}[label=3.\arabic*]
    \item \textbf{\llmsforresearcher}
    \begin{enumerate}[label=3.1.\arabic*]
        \item \annotators
        \item \judges
        \item \synthesis
        \item \subjects
    \end{enumerate}
    \item \textbf{\llmsforengineers}
    \begin{enumerate}[label=3.2.\arabic*]
        %\setcounter{enumii}{4}
        \item \llmusage
        \item \newtools
        \item \benchmarkingtasks
    \end{enumerate}
\end{enumerate}
\par\noindent\rule{\columnwidth}{0.5pt}

\subsection{LLMs as Tools for Software Engineering Researchers}
\label{sec:llms-as-tools-for-software-engineering-researchers}

LLMs can serve as powerful tools to help researchers conduct empirical studies.
They can automate various tasks such as data collection, pre-processing, and analysis.
For example, LLMs can apply predefined coding guides to qualitative datasets (\annotators), assess the quality of software artifacts (\judges), generate summaries of research papers (\synthesis), or simulate human behavior in empirical studies (\subjects).
This can significantly reduce the time and effort required to conduct a study.
However, besides the many advantages that all these applications bring, they also come with challenges such as potential threats to validity and implications for the reproducibility of study results.

%\par\noindent\hdashrule{\columnwidth}{0.5pt}{2pt 2pt}

\guidelinesubsubsection{LLMs as Annotators}
\label{sec:llms-as-annotators}

\paragraph{\bfseries Description}
Just like human annotators, LLMs can label artifacts based on a pre-defined coding guide.
However, they can label data much faster than any human could. 
In qualitative data analysis, manually annotating (``coding'') natural language text, e.g., in software artifacts, open-ended survey responses, or interview transcripts, is a time-consuming manual process~\cite{DBLP:journals/ase/BanoHZT24}.
LLMs can be used to augment or replace human annotations, provide suggestions for new codes (see Section \synthesis), or even automate the entire qualitative data analysis process.

\paragraph{\bfseries Example(s)}
Recent work in software engineering has begun to explore the use of LLMs for annotation tasks.
\citeauthor{Huang2023Enhancing}~\cite{Huang2023Enhancing} proposed an approach that utilizes multiple LLMs for joint annotation of mobile application reviews. 
They used three models of comparable size with an absolute majority voting rule (i.e., a label is only accepted if it receives more than half of the total votes from the models).
Accordingly, the annotations fell into three categories: exact matches (where all models agreed), partial matches (where a majority agreed), and non-matches (where no majority was reached).
The study by \citeauthor{DBLP:conf/msr/AhmedDTP25}~\cite{DBLP:conf/msr/AhmedDTP25} examined LLMs as annotators in software engineering research across five datasets, six LLMs, and ten annotation tasks.
They found that model-model agreement strongly correlates with human-model agreement, suggesting situations in which LLMs could effectively replace human annotators. Their research showed that for tasks where humans themselves disagreed significantly , models also performed poorly.
Conversely, if multiple LLMs reach similar solutions independently, then LLMs are likely suitable for the annotation task.
They proposed to use model confidence scores to identify specific samples that could be safely delegated to LLMs, potentially reducing human annotation effort without compromising inter-rater agreement.

\paragraph{\bfseries Advantages}
Recent research demonstrates several advantages of using LLMs as annotators, including their cost effectiveness and accuracy.
LLM-based annotation can dramatically reduce costs compared to human labeling, with studies showing cost reductions of 50-96\% on various natural language tasks~\cite{DBLP:conf/emnlp/WangLXZZ21}.
For example, \citeauthor{DBLP:conf/chi/HeHDRH24} found in their study that GPT-4 annotation costs only \$122.08 compared to \$4,508 for a comparable MTurk pipeline~\cite{DBLP:conf/chi/HeHDRH24}.
Moreover, the LLM-based approach resulted in a completion time of just 2 days versus several weeks for the crowd-sourced approach.
LLMs consistently demonstrate strong performance, with ChatGPT's accuracy exceeding crowd workers by approximately 25\% on average~\cite{DBLP:journals/corr/abs-2303-15056}, achieving impressive results in specific tasks such as sentiment analysis (65\% accuracy)~\cite{DBLP:journals/corr/abs-2304-10145}.
LLMs also show remarkably high inter-rater agreement, higher than crowd workers and trained annotators~\cite{DBLP:journals/corr/abs-2303-15056}.

\paragraph{\bfseries Challenges}
Challenges of using LLMs as annotators include reliability issues, human-LLM interaction challenges, biases, errors, and resource considerations.
Studies suggest that while LLMs show promise as annotation tools in SE research, their optimal use may be in augmenting rather than replacing human annotators~\cite{DBLP:conf/emnlp/WangLXZZ21, DBLP:conf/chi/HeHDRH24}.
LLMs can negatively affect human judgment when labels are incorrect~\cite{DBLP:conf/www/HuangKA23a}, and their overconfidence requires careful verification~\cite{DBLP:conf/kdd/WanSJKCNSSWYABJ24}.
Moreover, in previous studies, LLMs have shown significant variability in annotation quality depending on the dataset and the annotation task~\cite{DBLP:journals/corr/abs-2306-00176}.
Studies have shown that LLMs are especially unreliable for high-stakes labeling tasks~\cite{DBLP:conf/chi/Wang0RMM24} and that LLMs can have notable performance disparities between label categories~\cite{DBLP:journals/corr/abs-2304-10145}.
Recent empirical evidence indicates that LLM consistency in text annotation often falls below scientific reliability thresholds, with outputs being sensitive to minor prompt variations~\cite{DBLP:journals/corr/abs-2304-11085}.
Context-dependent annotations pose a specific challenge, as LLMs show difficulty in correctly interpreting text segments that require broader contextual understanding~\cite{DBLP:conf/chi/HeHDRH24}.
Although pooling of multiple outputs can improve reliability, this approach requires additional computational resources and still requires validation against human-annotated data.
While generally cost-effective, LLM annotation requires careful management of per token charges, particularly for longer texts~\cite{DBLP:conf/emnlp/WangLXZZ21}. Furthermore, achieving reliable annotations may require multiple runs of the same input to enable majority voting~\cite{DBLP:journals/corr/abs-2304-11085}, although the exact cost comparison between LLM-based and human annotation is controversial~\cite{DBLP:conf/chi/HeHDRH24}.
Finally, research has identified consistent biases in label assignment, including tendencies to overestimate certain labels and misclassify neutral content~\cite{DBLP:journals/corr/abs-2304-10145}.

%\par\noindent\hdashrule{\columnwidth}{0.5pt}{2pt 2pt}

\guidelinesubsubsection{LLMs as Judges}
\label{sec:llms-as-judges}

\paragraph{\bfseries Description}

LLMs can act as judges or raters to evaluate properties of software artifacts.
For example, LLMs can be used to assess code readability, adherence to coding standards, or the quality of code comments.
Judgment is distinct from the more qualitative task of assigning a code or label to typically unstructured text (see Section \annotators).
It is also different from using LLMs for general software engineering tasks, as discussed in Section \llmsforresearcher.

\paragraph{\bfseries Example(s)}

\citeauthor{DBLP:conf/re/LubosFTGMEL24}~\cite{DBLP:conf/re/LubosFTGMEL24} leveraged \href{https://www.llama.com/llama2/}{Llama-2} to evaluate the quality of software requirements statements. 
They prompted the LLM with the text below, where the words in curly brackets reflect the study parameters:

\begin{plain}
Your task is to evaluate the quality of a
software requirement.
Evaluate whether the following requirement is
{quality_characteristic}.
{quality_characteristic} means:
{quality_characteristic_explanation}
The evaluation result must be: 'yes' or 'no'.
Request: Based on the following description of
the project:
{project_description}
Evaluate the quality of the following
requirement: {requirement}.
Explain your decision and suggest an improved version.
\end{plain}

\citeauthor{DBLP:conf/re/LubosFTGMEL24} evaluated the LLM's output against human judges, to assess how well the LLM matched experts. 
Agreement can be measured in many ways; this study used Cohen's kappa~\cite{cohen60} and found moderate agreement for simple requirements and poor agreement for more complex requirements.
However, human evaluation of machine judgments may not be the same as evaluation of human judgments and is an open area of research~\cite{DBLP:journals/corr/abs-2410-03775}. 
A crucial decision in studies focusing on LLMs as judges is the number of examples provided to the LLM.
This might involve no tuning (zero-shot), several examples (few-shot), or providing many examples (closer to traditional training data).
In the example above, \citeauthor{DBLP:conf/re/LubosFTGMEL24} chose zero-shot tuning, providing no specific guidance besides the provided context; they did not show the LLM what a `yes' or `no' answer might look like. 
In contrast, \citeauthor{wang2025multimodalrequirementsdatabasedacceptance}~\cite{wang2025multimodalrequirementsdatabasedacceptance} provided a rubric to an LLM when treating LLMs as a judge to produce interpretable scales (0 to 4) for the overall quality of acceptance criteria generated from user stories.
They combined this with a lower-level reward stage to refine acceptance criteria, which significantly improved correctness, clarity, and alignment with the user stories.

\paragraph{\bfseries Advantages}

Depending on the model configuration, LLMs can provide relatively consistent evaluations.
They can further help mitigate human biases and the general variability that human judges might introduce.
This may lead to more reliable and reproducible results in empirical studies, to the extent that these models can be reproduced or checkpointed.
LLMs can be much more efficient and scale more easily than the equivalent human approach.
With LLM automation, entire datasets can be assessed, as opposed to subsets, and the assessment produced can be at the selected level of granularity, e.g., binary outputs (`yes' or `no') or defined levels (e.g., 0 to 4).
However, the main constraint that varies by model and budget is the input context size, i.e., the number of tokens one can pass into a model.
For example, the upper bound of the context that can be passed to OpenAi's \textsf{o1-mini} model is \href{https://help.openai.com/en/articles/9855712-openai-o1-models-faq-chatgpt-enterprise-and-edu}{32k tokens}.

\paragraph{\bfseries Challenges}

When relying on the judgment of LLMs, researchers must build a \textit{reliable} process for generating judgment labels that considers the non-deterministic nature of LLMs and report the intricacies of that process transparently~\cite{DBLP:journals/corr/abs-2412-12509}.
For example, the order of options has been shown to affect LLM outputs in multiple-choice settings~\cite{DBLP:conf/naacl/PezeshkpourH24}. 
In addition to reliability, other quality attributes to consider include the \textit{accuracy} of the labels. % and the speed and \textit{scalability} of the process. 
For example, a reliable LLM might be reliably inaccurate and wrong. 
Evaluating and judging large numbers of items, for example, to perform fault localization on the thousands of bugs that large open-source projects have to deal with, comes with costs in clock time, compute time, and environmental sustainability.
Evidence shows that LLMs can behave differently when reviewing their own outputs~\cite{NEURIPS2024_7f1f0218}.
In more human-oriented datasets (such as discussions of pull requests) LLMs may suffer from well-documented biases and issues with fairness~\cite{Gallegos2024BiasAF}. 
For tasks in which human judges disagree significantly, it is not clear if an LLM judge should reflect the majority opinion or act as an independent judge.
The underlying statistical framework of an LLM usually pushes outputs towards the most likely (majority) answer. 
There is ongoing research on how suitable LLMs are as independent judges.
Questions about bias, accuracy, and trust remain~\cite{DBLP:journals/corr/abs-2406-18403}.
There is reason for concern about LLMs judging student assignments or doing peer review of scientific papers~\cite{DBLP:conf/coling/ZhouC024}.
Even beyond questions of technical capacity, ethical questions remain, particularly if there is some implicit expectation that a human is judging the output.
Involving a human in the judgment loop, for example, to contextualize the scoring, is one approach~\cite{panHumanCenteredDesignRecommendations2024}. 
However, the lack of large-scale ground truth datasets for benchmarking LLM performance in judgment studies hinders progress in this area.

%\par\noindent\hdashrule{\columnwidth}{0.5pt}{2pt 2pt}

\guidelinesubsubsection{LLMs for Synthesis}
\label{sec:llms-for-synthesis}

\paragraph{\bfseries Description}

LLMs can support synthesis tasks in software engineering research by processing and distilling information from qualitative data sources.
In this context, synthesis refers to the process of integrating and interpreting information from multiple sources to generate higher-level insights, identify patterns across datasets, and develop conceptual frameworks or theories.
Unlike annotation (see Section \annotators), which focuses on categorizing or labeling individual data points, synthesis involves connecting and interpreting these annotations to develop a cohesive understanding of the phenomenon being studied.
Synthesis tasks can also mean generating synthetic datasets (e.g., source code, bug-fix pairs, requirements, etc.) that are then used in downstream tasks to train, fine-tune, or evaluate existing models or tools.
In this case, the synthesis is done primarily using the LLM and its training data; the input is limited to basic instructions and examples.

\paragraph{\bfseries Example(s)} 

Published examples of applying LLMs for synthesis in the software engineering domain are still scarce.
However, recent work has explored the use of LLMs for qualitative synthesis in other domains, allowing reflection on how LLMs can be applied for this purpose in software engineering~\cite{DBLP:journals/ase/BanoHZT24}.
\citeauthor{barros2024largelanguagemodelqualitative}~\cite{barros2024largelanguagemodelqualitative} conducted a systematic mapping study on the use of LLMs for qualitative research.
They identified examples in domains such as healthcare and social sciences (see, e.g., \cite{de2024performing,mathis2024inductive}) in which LLMs were used to support qualitative analysis, including grounded theory and thematic analysis.
Overall, the findings highlight the successful generation of preliminary coding schemes from interview transcripts, later refined by human researchers, along with support for pattern identification.
This approach was reported not only to expedite the initial coding process but also to allow researchers to focus more on higher-level analysis and interpretation.
However, \citeauthor{barros2024largelanguagemodelqualitative} emphasize that effective use of LLMs requires structured prompts and careful human oversight.
%This particular paper suggests using LLMs to support tasks such as initial coding and theme identification while conservatively reserving interpretative or creative processes for human analysts.
Similarly, \citeauthor{leça2024applicationsimplicationslargelanguage}~\cite{leça2024applicationsimplicationslargelanguage} conducted a systematic mapping study to investigate how LLMs are used in qualitative analysis and how they can be applied in software engineering research.
Consistent with the study by \citeauthor{barros2024largelanguagemodelqualitative}~\cite{barros2024largelanguagemodelqualitative}, \citeauthor{leça2024applicationsimplicationslargelanguage} identified that LLMs are applied primarily in tasks such as coding, thematic analysis, and data categorization, reducing the time, cognitive demands, and resources required for these processes.
Finally, \citeauthor{DBLP:journals/corr/abs-2506-21138}'s work is an example of using LLMs to create synthetic datasets.
They present an approach to generate synthetic requirements, showing that they \enq{can match or surpass human-authored requirements for specific classification tasks}~\cite{DBLP:journals/corr/abs-2506-21138}.

\paragraph{\bfseries Advantages}

LLMs offer promising support for synthesis in SE research by helping researchers process artifacts such as interview transcripts and survey responses, or by assisting literature reviews.
Qualitative research in SE traditionally faces challenges such as limited scalability, inconsistencies in coding, difficulties in generalizing findings from small or context-specific samples, and the influence of the researchers' subjectivity on data interpretation~\cite{DBLP:journals/ase/BanoHZT24}. 
The use of LLMs for synthesis can offer advantages in addressing these challenges~\cite{DBLP:journals/ase/BanoHZT24, barros2024largelanguagemodelqualitative, leça2024applicationsimplicationslargelanguage}.
LLMs can reduce manual effort and subjectivity, improve consistency and generalizability, and assist researchers in deriving codes and developing coding guides during the early stages of qualitative data analysis~\cite{DBLP:conf/chi/ByunVS23,DBLP:journals/ase/BanoHZT24}.
LLMs can enable researchers to analyze larger datasets, identifying patterns across broader contexts than traditional qualitative methods typically allow. In addition, they can help mitigate the effects of human subjectivity.
However, while LLMs streamline many aspects of qualitative synthesis, careful oversight remains essential.
% Models can develop coding guides and identify emerging themes, potentially automating the entire synthesis process across large qualitative datasets.

\paragraph{\bfseries Challenges}

Although LLMs have the potential to automate synthesis, concerns about overreliance remain, especially due to discrepancies between AI- and human-generated insights, particularly in capturing contextual nuances~\cite{bano2023exploringqualitativeresearchusing}.
\citeauthor{bano2023exploringqualitativeresearchusing}~\cite{bano2023exploringqualitativeresearchusing} found that while LLMs can provide structured summaries and qualitative coding frameworks, they may misinterpret nuanced qualitative data due to a lack of contextual understanding.
Other studies have echoed similar concerns~\cite{DBLP:journals/ase/BanoHZT24, barros2024largelanguagemodelqualitative, leça2024applicationsimplicationslargelanguage}.
In particular, LLMs cannot independently assess the validity of arguments.
Critical thinking remains a human responsibility in qualitative synthesis.
\citeauthor{peterschineyee2025generalizationbias} have shown that popular LLMs and LLM-based tools tend to overgeneralize results when summarizing scientific articles, that is, they \enq{produce broader generalizations of scientific results than those in the original.}
Compared to human-authored summaries, \enq{LLM summaries were nearly five times more likely to contain broad generalizations}~\cite{peterschineyee2025generalizationbias}.
In addition, LLMs can produce biased results, reinforcing existing prejudices or omitting essential perspectives, making human oversight crucial to ensure accurate interpretation, mitigate biases, and maintain quality control.
Moreover, the proprietary nature of many LLMs limits transparency.
In particular, it is unknown how the training data might affect the synthesis process.
Furthermore, reproducibility issues persist due to the influence of model versions and prompt variations on the synthesis result.

%\par\noindent\hdashrule{\columnwidth}{0.5pt}{2pt 2pt}

% Reverted back to "subject", as participant seems to be primarily used for humans: see  https://dictionary.apa.org/subject and https://dictionary.apa.org/participant
\guidelinesubsubsection{LLMs as Subjects}
\label{sec:llms-as-subjects}

\paragraph{\bfseries Description}

In empirical studies, data is collected from participants through methods such as surveys, interviews, or controlled experiments.
LLMs can serve as \emph{subjects} in empirical studies by simulating human behavior and interactions (we use the term ``subject'' since ``participant'' implies being human~\cite{apa-dict-subject}).
In this capacity, LLMs generate responses that approximate those of human participants, which makes them particularly valuable for research involving user interactions, collaborative coding environments, and software usability assessments.
This approach enables data collection that closely reflects human reactions while avoiding the need for direct human involvement.
To achieve this, prompt engineering techniques are widely employed, with a common approach being the use of the \textit{Personas Pattern}~\cite{DBLP:journals/corr/abs-2308-07702}, which involves tailoring LLM responses to align with predefined profiles or roles that emulate specific user archetypes.
\citeauthor{ZHAO2025101167} outline opportunities and challenges of using LLMs as research subjects in detail~\cite{ZHAO2025101167}.
Furthermore, recent sociological studies have emphasized that, to be effectively utilized in this capacity, LLMs---including their agentic versions---should meet four criteria of algorithmic fidelity~\cite{DBLP:journals/corr/abs-2209-06899}.
Generated responses should be: (1) indistinguishable from human-produced texts (e.g., LLM-generated code reviews should be comparable to those from real developers); (2) consistent with the attitudes and sociodemographic information of the conditioning context (e.g., LLMs simulating junior developers should exhibit different confidence levels, vocabulary, and concerns compared to senior engineers); (3) naturally aligned with the form, tone, and content of the provided context (e.g., responses in an agile stand-up meeting simulation should be concise, task-focused, and aligned with sprint objectives rather than long, formal explanations); and (4) reflective of patterns in relationships between ideas, demographics, and behavior observed in comparable human data (e.g., discussions on software architecture decisions should capture trade-offs typically debated by human developers, such as maintainability versus performance, rather than abstract theoretical arguments).

\paragraph{\bfseries Example(s)}

LLMs can be used as subjects in various types of empirical studies, enabling researchers to simulate human participants.
The broader applicability of LLM-based studies beyond software engineering has been compiled by \citeauthor{DBLP:journals/ipm/XuSRGPLSH24}~\cite{DBLP:journals/ipm/XuSRGPLSH24}, who examined various uses in social science research.
Given the socio-technical nature of software development, some of these approaches are transferable to empirical software engineering research.
For example, LLMs can be applied in survey and interview studies to impersonate developers responding to survey questionnaires or interviews, allowing researchers to test the clarity and effectiveness of survey items or to simulate responses under varying conditions, such as different levels of expertise or cultural contexts.
For example, \citeauthor{DBLP:journals/ase/GerosaTSS24}~\cite{DBLP:journals/ase/GerosaTSS24} explored persona-based interviews and multi-persona focus groups, demonstrating how LLMs can emulate human responses and behaviors while addressing ethical concerns, biases, and methodological challenges.
Another example are usability studies, in which LLMs can simulate end-user feedback, providing insights into potential usability issues and offering suggestions for improvement based on predefined user personas. This aligns with the work of \citeauthor{bano2025doessoftwareengineerlook}~\cite{bano2025doessoftwareengineerlook}, who investigated biases in LLM-generated candidate profiles in SE recruitment processes.
Their study, which analyzed both textual and visual inputs, revealed biases favoring male, Caucasian candidates, lighter skin tones, and slim physiques, particularly for senior roles.

\paragraph{\bfseries Advantages}

Using LLMs as subjects can reduce the effort of recruiting human participants, a process that is often time-consuming and costly~\cite{DBLP:conf/vl/Madampe0HO24}, by augmenting existing datasets or, in some cases, completely replacing human participants. 
In interviews and survey studies, LLMs can simulate diverse respondent profiles, enabling access to underrepresented populations, which can strengthen the generalizability of findings. 
In data mining studies, LLMs can generate synthetic data (see \synthesis) to fill gaps where real-world data is unavailable or underrepresented.

\paragraph{\bfseries Challenges}

It is important that researchers are aware of the inherent biases~\cite{Crowell2023} and limitations~\cite{DBLP:journals/ais/HardingDLL24, DBLP:journals/corr/abs-2402-01908} when using LLMs as study subjects. 
\citeauthor{schroeder2025llmspsychology}~\cite{schroeder2025llmspsychology}, who have studied discrepancies between LLM and human responses, even conclude that \enq{LLMs do not simulate human psychology and recommend that psychological researchers should treat LLMs as useful but fundamentally unreliable tools that need to be validated against human responses for every new application.}
One critical concern is construct validity.
LLMs have been shown to misrepresent demographic group perspectives, failing to capture the diversity of opinions and experiences within a group.
The use of identity-based prompts can reduce identities to fixed and innate characteristics, amplifying perceived differences between groups.
These biases introduce the risk that studies which rely on LLM-generated responses may inadvertently reinforce stereotypes or misrepresent real-world social dynamics.
Alternatives to demographic prompts can be employed when the goal is to broaden response coverage~\cite{DBLP:journals/corr/abs-2402-01908}.
%To mitigate these issues, encoded identity names, that is, statistically representative person names associated with specific demographic groups, can be used instead of explicit labels (e.g., using ``Mario Rossi'' instead of the label ``Italian white man'').
%Moreover, the temperature setting can be increased to enhance response diversity.
%All three strategies aim to reduce the risk of stereotyped, flattened, or overly deterministic outputs by introducing indirectness or variability into the prompt, thereby encouraging the model to generate responses that better reflect the internal diversity and nuance of human perspectives.
Beyond construct validity, internal validity must also be considered, particularly with regard to causal conclusions based on studies relying on LLM-simulated responses.
Finally, external validity remains a challenge, as findings based on LLMs may not generalize to humans.


\subsection{LLMs as Tools for Software Engineers}
\label{sec:llms-as-tools-for-software-engineers}

LLM-based assistants have become an essential tool for software engineers, supporting them in various tasks such as code generation and debugging.
Researchers have studied how software engineers use LLMs (\llmusage), developed new tools that integrate LLMs (\newtools), and benchmarked LLMs for software engineering tasks (\benchmarkingtasks).
%Like in the previous section, we outline the advantages that studying LLMs in these contexts brings, but also point to specific challenges.

%\par\noindent\hdashrule{\columnwidth}{0.5pt}{2pt 2pt}

\guidelinesubsubsection{Studying LLM Usage in Software Engineering}
\label{sec:studying-llm-usage-in-software-engineering}

\paragraph{\bfseries Description}

Studying how software engineers use LLMs is crucial to understand the current state of practice in SE.
Researchers can observe software engineers' usage of LLM-based tools in the field, or study if and how they adopt such tools, their usage patterns, as well as perceived benefits and challenges.
Surveys, interviews, observational studies, or analysis of usage logs can provide insights into how LLMs are integrated into development processes, how they influence decision making, and what factors affect their acceptance and effectiveness. 
Such studies can inform improvements for existing LLM-based tools, motivate the design of novel tools, or derive best practices for LLM-assisted software engineering.
They can also uncover risks or deficiencies of existing tools.

\paragraph{\bfseries Example(s)}

\citeauthor{DBLP:journals/pacmse/KhojahM0N24} investigated the use of ChatGPT (GPT-3.5) by professional software engineers in a week-long observational study~\cite{DBLP:journals/pacmse/KhojahM0N24}.
They found that most developers do not use the code generated by ChatGPT directly but instead use the output as a guide to implement their own solutions.
%Furthermore, they recommend that future research investigate the use of ChatGPT in non code-related SE tasks and for training and learning SE concepts.
\citeauthor{DBLP:conf/csee/AzanzaPIG24} conducted a case study that evaluated the impact of introducing LLMs on the onboarding process of new software developers~\cite{DBLP:conf/csee/AzanzaPIG24} (GPT-3).
Their study identified potential in the use of LLMs for onboarding, as it can allow newcomers to seek information on their own, without the need to \enq{bother} senior colleagues.
%However, they also highlight significant privacy concerns with the use of proprietary, third-party LLMs.
Surveys can help researchers to quickly provide an overview of the current perceptions of LLM usage.
For example, \citeauthor{DBLP:conf/icsa/JahicS24} surveyed participants from 15 software companies regarding their practices on LLMs in software engineering~\cite{DBLP:conf/icsa/JahicS24}.
They found that the majority of study participants had already adopted AI for software engineering tasks; most of them used ChatGPT.
Multiple participants cited copyright and privacy issues, as well as inconsistent or low-quality outputs, as barriers to adoption.
Retrospective studies that analyze data generated while developers use LLMs can provide additional insights into human-LLM interactions.
For example, researchers can employ data mining methods to build large-scale conversation datasets, such as the DevGPT dataset introduced by \citeauthor{DBLP:conf/msr/XiaoTHM24}~\cite{DBLP:conf/msr/XiaoTHM24}.
Conversations can then be analyzed using quantitative~\cite{DBLP:conf/msr/RabbiCZI24} and qualitative~\cite{DBLP:conf/msr/MohamedPP24} analysis methods.

\paragraph{\bfseries Advantages}

Studying the real-world usage of LLM-based tools allows researchers to understand the state of practice and guide future research directions.
In field studies, researchers can uncover usage patterns, adoption rates, and the influence of contextual factors on usage behavior.
Outside of a controlled study environment, researchers can uncover contextual information about LLM-assisted SE workflows beyond the specific LLMs being evaluated.
This may, for example, help researchers generate hypotheses about how LLMs impact developer productivity, collaboration, and decision-making processes.
In controlled laboratory studies, researchers can study specific phenomena related to LLM usage under carefully regulated, but potentially artificial conditions.
Specifically, they can isolate individual tasks in the software engineering workflow and investigate how LLM-based tools may support task completion.
Furthermore, controlled experiments allow for direct comparisons between different LLM-based tools.
The results can then be used to validate general hypotheses about human-LLM interactions in SE.

\paragraph{\bfseries Challenges}

When conducting field studies in real-world environments, researchers have to ensure that their study results are \enq{dependable}~\cite{Sullivan2011-ub} beyond the traditional validity criteria such as internal or construct validity.
The usage environment in a real-world context can often be extremely diverse.
The integration of LLMs can range from specific LLMs based on company policy to the unregulated use of any available LLM. 
Both extremes may influence the adoption of LLM by software engineers, and hence need to be addressed in the study methodology.
In addition, in longitudinal case studies, the timing of the study may have a significant impact on its result, as LLMs and LLM-based tools are rapidly evolving.
Moreover, developers are still learning how to best make use of the new technology, and best practices are still being developed and established.
Furthermore, the predominance of proprietary commercial LLM-based tools in the market poses a significant barrier to research.
Limited access to telemetry data or other usage metrics restricts the ability of researchers to conduct comprehensive analyses of real-world tool usage.
To make study results reliable, researchers must establish that their results are rigorous, for example, by using triangulation to understand and minimize potential biases~\cite{Sullivan2011-ub}.
When it comes to studying LLM usage in controlled laboratory studies, researchers may struggle with the inherent variability of LLM outputs.
Since reproducibility and precision are essential for hypothesis testing, the stochastic nature of LLM responses---where identical prompts may yield different outputs across participants---can complicate the interpretation of experimental results and affect the reliability of study findings.

%\par\noindent\hdashrule{\columnwidth}{0.5pt}{2pt 2pt}

\guidelinesubsubsection{LLMs for New Software Engineering Tools}
\label{sec:llms-for-new-software-engineering-tools}

\paragraph{\bfseries Description}

LLMs are being integrated into new tools that support software engineers in their daily tasks, e.g., to assist in code comprehension~\cite{DBLP:conf/chi/YanHWH24} and test case generation~\cite{DBLP:journals/tse/SchaferNET24}.
One way of integrating LLM-based tools into software engineers' workflows are GenAI agents.
Unlike traditional LLM-based tools, these agents are capable of acting autonomously and proactively, are often tailored to meet specific user needs, and can interact with external environments~\cite{takerngsaksiri2024human,wiesinger2025agents}.
From an architectural perspective, GenAI agents can be implemented in various ways~\cite{wiesinger2025agents}.
However, they generally share three key components: (1) a reasoning mechanism that guides the LLM (often enabled by advanced prompt engineering), (2) a set of tools to interact with external systems (e.g., APIs or databases), and (3) a user communication interface that extends beyond traditional chat-based interactions~\cite{DBLP:conf/icsm/RichardsW24, DBLP:journals/tmlr/SumersYN024, DBLP:journals/corr/abs-2309-07870}.
Researchers can also test and compare different tool architectures to increase artifact quality and developer satisfaction.

\paragraph{\bfseries Example(s)}

\citeauthor{DBLP:conf/chi/YanHWH24} proposed IVIE, a tool integrated into the VS Code graphical interface that generates and explains code using LLMs~\cite{DBLP:conf/chi/YanHWH24}.
The authors focused more on the presentation, providing a user-friendly interface to interact with the LLM. 
\citeauthor{DBLP:journals/tse/SchaferNET24}~\cite{DBLP:journals/tse/SchaferNET24} presented a large-scale empirical evaluation on the effectiveness of LLMs for automated unit test generation.
They presented TestPilot, a tool that implements an approach in which the LLM is provided with prompts that include the signature and implementation of a function under test, along with usage examples extracted from the documentation.
\citeauthor{DBLP:conf/icsm/RichardsW24} introduced a preliminary GenAI agent designed to assist developers in understanding source code by incorporating a reasoning component grounded in the theory of mind~\cite{DBLP:conf/icsm/RichardsW24}.

\paragraph{\bfseries Advantages}

From an engineering perspective, developing LLM-based tools is easier than implementing many traditional SE approaches such as static analysis or symbolic execution.
Depending on the capabilities of the underlying model, it is also easier to build tools that are independent of a specific programming language.
This enables researchers to build tools for a more diverse set of tasks.
In addition, it allows them to test their tools in a wider range of contexts.

\paragraph{\bfseries Challenges}

Traditional approaches, such as static analysis, are deterministic. LLMs are not.
Although the non-determinism of LLMs can be mitigated using configuration parameters and prompting strategies, this poses a major challenge.
It can be challenging for researchers to evaluate the effectiveness of a tool, as minor changes in the input can lead to major differences in the performance.
Since the exact training data is often not published by model vendors, a reliable assessment of tool performance for unknown data difficult.
From an engineering perspective, while open models are available, the most capable ones require substantial hardware resources.
Using cloud-based APIs or relying on third-party providers for hosting, while seemingly a potential solution, introduces new concerns related to data privacy and security.

%\par\noindent\hdashrule{\columnwidth}{0.5pt}{2pt 2pt}

\guidelinesubsubsection{Benchmarking LLMs for Software Engineering Tasks}
\label{sec:benchmarking-llms-for-software-engineering-tasks}

\paragraph{\bfseries Description}

Benchmarking is the process of evaluating an LLM's performance using standardized tasks and metrics, which requires high-quality reference datasets.
LLM output is compared to a ground truth from the benchmark dataset using general metrics for text generation, such as \emph{ROUGE}, \emph{BLEU}, or \emph{METEOR}~\cite{10.1145/3695988}, or task-specific metrics, such as \emph{CodeBLEU} for code generation.
For example, \emph{HumanEval}~\cite{DBLP:journals/corr/abs-2107-03374} is often used to assess code generation, establishing it as a de facto standard.

\paragraph{\bfseries Example(s)}

In SE, benchmarking may include the evaluation of an LLM's ability to produce accurate and reliable outputs for a given input, usually a task description, which may be accompanied by data obtained from curated real-world projects or from synthetic SE-specific datasets.
Typical tasks include code generation, code summarization, code completion, and code repair, but also natural language processing tasks, such as anaphora resolution (i.e., the task of identifying the referring expression of a word or phrase occurring earlier in the text). %, interesting for subfields such a Requirements Engineering. 
\emph{RepairBench}~\cite{silva2024repairbench}, for example, contains 574 buggy Java methods and their corresponding fixed versions, which can be used to evaluate the performance of LLMs in code repair tasks.
This benchmark uses the \emph{Plausible@1} metric (i.e., the probability that the first generated patch passes all test cases) and the \emph{AST Match@1} metric (i.e., the probability that the abstract syntax tree of the first generated patch matches the ground truth patch).
\emph{SWE-Bench}~\cite{DBLP:conf/iclr/JimenezYWYPPN24} is a more generic benchmark that contains 2,294 SE Python tasks extracted from GitHub pull requests.
To score the LLM's performance on the tasks, the benchmark validates whether the generated patch is applicable (i.e., successfully compiles) and calculates the percentage of passed test cases.

\paragraph{\bfseries Advantages}

Properly built benchmarks provide objective, reproducible evaluation across different tasks, enabling a comparison between different models (and versions).
In addition, benchmarks built for specific SE tasks can help identify LLM weaknesses and support their optimization and fine-tuning for such tasks.
They can foster open science practices by providing a common ground for sharing data (e.g., as part of the benchmark itself) and results (e.g., of models run against a benchmark).
Benchmarks built using real-world data can help legitimize research results for practitioners, supporting industry-academia collaboration.
However, benchmarks are subject to several challenges.

\paragraph{\bfseries Challenges}

Benchmark contamination, that is, the inclusion of the benchmark in the LLM training data~\cite{DBLP:journals/corr/abs-2410-16186}, has recently been identified as a problem.
The careful selection of samples and the creation of the corresponding input prompts is particularly important, as correlations between prompts may bias the benchmark results~\cite{DBLP:conf/acl/SiskaMAB24}.
Although LLMs might perform well on a specific benchmark such as \emph{HumanEval} or \emph{SWE-bench}, they do not necessarily perform well on other benchmarks.
Moreover, benchmarks usually do not capture the full complexity and diversity of software engineering work~\cite{Chandra2025benchmarks}.
Another issue is that metrics used for benchmarks, such as \emph{perplexity} or \emph{BLEU-N}, do not necessarily reflect human judgment.
Recently, \citeauthor{cao2025should}~\cite{cao2025should} have proposed guidelines for the creation of LLM benchmarks related to coding tasks, grounded in a systematic survey of existing benchmarks. 
In this process, they highlight current shortcomings related to reliability, transparency, irreproducibility, low data quality, and inadequate validation measures.
For more details on benchmarks, see Section \benchmarksmetrics. 


\section{Guidelines}
\label{sec:guidelines}

Our guidelines focus on LLMs, that is, foundational models that use \emph{text} as in- and output.
We do not address multi-modal foundation models that support other media such as images.
However, many of our recommendations apply to these models as well.
We consider the extension of our guidelines to explicitly address multi-modal foundational models an important direction for future work.

The main goal of our guidelines is to \emph{enable reproducibility and replicability} of empirical studies involving LLMs in software engineering.
While we consider LLM-based studies to have characteristics that differ from traditional empirical studies (e.g., their inherent non-determinism and the fact that truly open models are rare), previous guidelines regarding open science and empirical studies still apply.
Although full reproducibility of LLM study results is very challenging given LLM's non-determinism, transparency on LLM usage, methods, data, and architecture, as suggested by our guidelines, is an essential prerequisite for future replication studies.

The wording of our guidelines (\must, \should, \may) follows RFC 2119~\cite{rfc2119} and RFC 8174~\cite{rfc8174}.
Throughout the following sections, we mention the information we expect researchers to report in the \paper and/or in the \supplementarymaterial.
We are aware that different publication venues have different page limits and that not all aspects can be reported in the \paper.
If information \must be reported in the \paper, we explicitly mention this in the specific guidelines.
Of course, it is better to report essential information that we expect in the \paper in the \supplementarymaterial than not at all.
The \supplementarymaterial \should be published according to the \emph{ACM SIGSOFT Open Science Policies}~\cite{daniel_graziotin_2024_10796477}.
%Unless explicitly mentioned otherwise, we expect information to be reported either in the \paper and/or the \supplementarymaterial.

At the beginning of each section, we provide a brief \underline{\emph{tl;dr}} summary that lists the most important aspects of the corresponding guideline.
In addition to our \textbf{recommendations}, we provide \textbf{examples} from the SE research community and beyond, as well as the \textbf{advantages} and potential \textbf{challenges} of following the respective guidelines.
We conclude each guideline by linking it to the \textbf{study types}.
The remainder of this section presents our eight guidelines and is structured as follows:

\par\noindent\rule{\columnwidth}{0.5pt}
\begin{enumerate}[label=4.\arabic*]
    \item \usagerole
    \item \modelversion
    \item \toolarchitecture
    \item \prompts
    \item \humanvalidation
    \item \openllm
    \item \benchmarksmetrics
    \item \limitationsmitigations
\end{enumerate}
\par\noindent\rule{\columnwidth}{0.5pt}


\subsection{Declare LLM Usage and Role}
\label{sec:declare-llm-usage-and-role}

\vspace{0.5\baselineskip}
\begin{framed}
\underline{tl;dr:} Researchers \must disclose any use of LLMs to support empirical studies in their \paper. They \should report why they used an LLM, what tasks were automated tasks, and the expected benefits of this automation.
\end{framed}

\guidelinesubsubsection{Rationale}

Transparency about LLM involvement is a prerequisite for informed assessment of a study's scope, limitations, and potential biases.
Without explicit disclosure, readers cannot evaluate how the LLM's characteristics may have influenced the research process or its outcomes.

\guidelinesubsubsection{Recommendations}

When conducting any kind of empirical study involving LLMs, researchers \must clearly declare that an LLM was used (see \scope for what we consider relevant research support).
This \should be done in a suitable section of the \paper, for example, in the introduction or research methods section; \citeauthor{cheng2025writing}~\cite{cheng2025writing} argue specifically for the methods section, since acknowledgments are easily missed at the end of a paper.
% GROUNDING cheng2025writing: "the most helpful and transparent means for describing the use of these tools is within the methods section of a manuscript"
The ACM Policy on Authorship requires authors to describe in the methods section any use of AI in the research itself~\cite{ACM2026}.
% GROUNDING ACM2026: "the specific use(s) of AI tools must be described in detail in the methods section of the Work"

Beyond generic declarations, researchers \should report the exact purpose of using an LLM in a study, the tasks it was used to automate, and the expected benefits in the \paper.
A sufficient declaration specifies not only \emph{that} an LLM was used, but also \emph{which} LLM (name and version), \emph{how} it was used (e.g., as an annotator, code generator, or judge), and \emph{where} in the research process it was employed (e.g., data collection, analysis, or synthesis).

When the LLM is central to the study (e.g., as the main tool being evaluated or as a core component of the research method), the declaration \should be prominent and detailed, appearing in the methodology section with cross-references to the specific guidelines that apply (e.g., Sections \modelversion, \design, and \traces).
When the LLM's role is more tangential (e.g., used for a single preprocessing step), a brief but explicit statement in the methodology section is sufficient.
When a study assigns multiple distinct roles to LLMs (e.g., one model generates evaluation data while another scores outputs), each role \should be declared separately.
In each case, the disclosure \must be specific enough for readers to assess how the LLM's involvement may affect the study's validity and reproducibility.

\guidelinesubsubsection{Examples}

The \emph{ACM Policy on Authorship} requires reporting AI-generated artifacts such as code, datasets, and figures where they underlie a study's conclusions~\cite{ACM2026}.
% GROUNDING ACM2026: "the creation of artifacts that are directly relevant to the conclusions of the research, such as code, datasets, and charts or figures that rely on the AI tools"
Reporting in the methods section keeps this disclosure visible under double-blind review, where end-of-paper acknowledgments are often removed.
An example of an LLM disclosure beyond writing support can be found in a recent paper by \citeauthor{DBLP:conf/re/LubosFTGMEL24}~\cite{DBLP:conf/re/LubosFTGMEL24}, in which they write in the methodology section:
% GROUNDING DBLP:conf/re/LubosFTGMEL24: "We conducted an LLM-based evaluation of requirements utilizing the Llama 2 language model with 70 billion parameters, fine-tuned to complete chat responses."

\begin{quote}
\it
``We conducted an LLM-based evaluation of requirements utilizing the Llama 2 language model with 70 billion parameters, fine-tuned to complete chat responses...''
\end{quote}

\noindent A more contemporary declaration could similarly state:

\begin{quote}
\it
``We used Claude Opus 4.7 via the Anthropic API to synthesize themes from interview transcripts, with all prompts and conversation logs published as supplementary material.''
\end{quote}

\citeauthor{DBLP:journals/corr/abs-2601-11895}'s DevBench paper illustrates separate disclosure of multiple LLM roles: \emph{GPT-4o} generates the synthetic benchmark instances, nine models (including \emph{Claude 4 Sonnet}, \emph{GPT-4.1}, \emph{DeepSeek-V3}, and \emph{Ministral-3B}) are the evaluation subjects, and \emph{o3-mini} is the LLM judge that scores completions for relevance and helpfulness~\cite{DBLP:journals/corr/abs-2601-11895}.
% GROUNDING DBLP:journals/corr/abs-2601-11895: "Synthetic instances were generated with OpenAI's GPT-4o, chosen for its fluency, reasoning, and code generation capabilities"

\guidelinesubsubsection{Benefits}

Transparency in the use of LLMs helps other researchers understand the context and scope of the study, supporting interpretation and comparison of the results.
Realizing these benefits requires reporting the LLM's exact role and version (see \modelversion).

\guidelinesubsubsection{Challenges}

Declaring LLM usage requires only a brief statement and no additional experiments, making compliance straightforward.
One challenge might be authors' reluctance to disclose LLM usage for valid use cases, because they fear that AI-generated content makes reviewers think that the authors' work is less original.
In fact, there is evidence suggesting that AI disclosure can negatively affect trust in authors~\cite{SCHILKE2025104405}.
% GROUNDING SCHILKE2025104405: "Thirteen experiments consistently demonstrate that actors who disclose their AI usage are trusted less than those who do not."
However, the \emph{ACM Policy on Authorship} requires disclosure of AI used in the research itself, not AI used only to assist with writing~\cite{ACM2026}.
% GROUNDING ACM2026: "ACM no longer requires the disclosure of information regarding the use of AI (as distinct from AI used in the conduct of the research itself"
Our guidelines focus on such use (see \scope), not on proofreading or writing support.

\guidelinesubsubsection{Study Types}

Researchers \must follow this guideline for all study types.
The specific focus of the declaration varies by study type.
For \annotators, \judges, \synthesis, and \subjects, researchers \must declare the specific role assigned to the LLM (e.g., annotator, judge, synthesizer, or simulated participant).
For \llmusage, researchers \must clarify which LLM(s) the observed participants used and under which conditions.
For \newtools, researchers \must declare the LLM's role within the tool architecture and its contribution to the tool's functionality.
For \benchmarkingtasks, researchers \must declare which LLMs were benchmarked and for which tasks.

\guidelinesubsubsection{Advice for Reviewers}

The most common problem with disclosure is incompleteness or vagueness about how the LLM was used. If the paper says ``we used LLM X to help with task Y'' without specifying how, reviewers should request clarification. Such requests are typically minor revisions unless the missing details may reveal methodological problems.

If undisclosed LLM use is suspected, the reviewer should consult their editor or program chair. When the evidence is conclusive, the key question is the degree to which undisclosed use affects the study's contribution, ranging from negligible (e.g., word choice in a single sentence) to severe (e.g., generating the reported data or results).

\guidelinesubsubsection{See Also}
\begin{itemize}[label=$\rightarrow$]
    \item \seemodelversion: The disclosure is incomplete without naming the specific model and version.
    \item \seedesign: When the LLM lives inside a tool or agent, authors must also describe that tool's architecture and prompts.
    \item \seetraces: Session traces show what the LLM did during the study.
\end{itemize}


\subsection{Report Model Version, Configuration, and Customizations}
\label{sec:report-model-version-configuration-and-customizations}

\vspace{0.5\baselineskip}
\begin{framed}
\underline{tl;dr:} Researchers \must report the exact LLM model or tool version, configuration, and experiment date in the \paper. When using quantized models, researchers \should report the quantization level and method. For fine-tuned models, they \must describe the fine-tuning goal, dataset, and procedure in the \paper. Researchers \should include default parameters, explain model choices, compare base- and fine-tuned model using suitable metrics and benchmarks, and share fine-tuning data and weights as \supplementarymaterial (or alternatively justify in the \paper why they cannot share them). 
\end{framed}

\input{_guidelines/02_report-model-version-configuration-and-customizations.tex}


\subsection{Report Tool Architecture Beyond Models}
\label{sec:report-tool-architecture-beyond-models}

\vspace{0.5\baselineskip}
\begin{framed}
\underline{tl;dr:} \input{_tldr/03_report-tool-architecture-beyond-models-tldr.tex} 
\end{framed}

\input{_guidelines/03_report-tool-architecture-beyond-models.tex}


\subsection{Report Prompts, their Development, and Interaction Logs}
\label{sec:report-prompts-their-development-and-interaction-logs}

\vspace{0.5\baselineskip}
\begin{framed}
\underline{tl;dr:} \input{_tldr/04_report-prompts-their-development-and-interaction-logs-tldr.tex}
\end{framed}

\input{_guidelines/04_report-prompts-their-development-and-interaction-logs.tex}


\subsection{Use Human Validation for LLM Outputs}
\label{sec:use-human-validation-for-llm-outputs}

\vspace{0.5\baselineskip}
\begin{framed}
\underline{tl;dr:} \input{_tldr/05_use-human-validation-for-llm-outputs-tldr.tex}  
\end{framed}

\guidelinesubsubsection{Recommendations}

When LLMs are used to automate research or software development tasks previously performed by humans, it is essential to assess the LLM's performance. When existing reference datasets or comparison metrics cannot fully capture the target qualities relevant in the study context, researchers \should rely on human judgment to validate the LLM outputs.

Integrating human participants in the study design will require additional considerations, including a recruitment strategy, annotation guidelines, training sessions, or ethical approvals.
Therefore, researchers \should consider human validation early in the study design, not as an afterthought.
Authors \must clearly define the constructs that the human and LLM annotators evaluate~\cite{DBLP:conf/ease/RalphT18}.
When designing custom instruments to assess LLM output (e.g., questionnaires, scales), researchers \should share their instruments in the \supplementarymaterial.

When the judgment is subjective (i.e. depends on the judge's values or implicit or explicit theories of the phenomenon in question), the LLM output \should be evaluated against judgments aggregated from multiple human judges. In these cases, researchers \should clearly describe their aggregation method and document their reasoning. \added{The preferred method is to randomly order the objects requiring judgment and put them in groups small enough to judge in 1--3 hours. Two or three human experts review the first group together, simultaneously judging, discussing, and creating a set of decision rules documenting their reasoning. Then, the experts iterate among rounds of:}
\begin{itemize}
    \item independently rating one group of objects
    \item calculating inter-rater agreement (IRA) or reliability (IRR)
    \item meeting to discuss disagreements and reach consensus
    \item updating the decision rules so the disagreement will not recur.
\end{itemize}

\added{The goal of this process is to reach some target IRA or IRR (e.g. Krippendorff's $\alpha>0.8$), which indicates that the decision rules are sufficient. If the number of judgments required is large, it is permissible to continue with a single rater per judgment once the IRA/IRR target is reached and sustained for 2--3 groups.} 

Researchers \should provide measures of IRA or IRR, preferably broken down by round. Confounding factors \should be controlled for (e.g., by categorizing participants according to their level of experience or expertise).
Where applicable, researchers \should perform a power analysis to estimate the required sample size, ensuring sufficient statistical power in their experimental design.

Researchers \should use established reference models to compare humans with LLMs.
For example, \citeauthor{Schneider2025ReferenceModel}~\cite{Schneider2025ReferenceModel} outline design considerations for studies comparing LLMs with humans.
%They outline the relevant groups (researchers, experimenters, subjects), activities (study planning, task assignment, results analysis, interpretation), and artifacts (study design, tasks, results, conclusion).
%Further, they illustrate the relationship between these entities, the order of steps, and which design decisions need to be considered.
If studies aim to automate annotating software artifacts using an LLM, researchers \should follow systematic approaches to decide whether and how human annotators can be replaced (e.g. showing that a jury of three LLMs exhibit model-to-model agreement exceeding some threshold).  
% For example, \citeauthor{DBLP:conf/msr/AhmedDTP25}~\cite{DBLP:conf/msr/AhmedDTP25} suggest a method that involves using a jury of three LLMs with 3 to 4 few-shot examples rated by humans, where the model-to-model agreement on all samples is determined using \emph{Krippendorff's Alpha}.
% If the agreement is high (alpha > 0.5), a human rating can be replaced with an LLM-generated one.
% In cases of low model-to-model agreement (alpha $\le$ 0.5), they then evaluate the prediction confidence of the model, selectively replacing annotations where the model confidence is high ($\ge$ 0.8).

Besides generating and modifying content, agentic software development tools such as Claude Code can autonomously call command-line tools or pull in additional information from MCP servers.
For such tools, it makes sense to separate textual output (e.g., file changes) from output related to tool execution.
When evaluating agentic tools, researchers \should assess the feedback that users provided for proposed changes, report statistics on how frequently they accepted the content, and how they modified it.
This agentic human-in-the-loop interaction approach is essentially a built-in human validation, even though the degrees of freedom are larger than in more traditional experiments that use LLMs directly.

\guidelinesubsubsection{Example(s)}

\citeauthor{DBLP:conf/msr/AhmedDTP25}~\cite{DBLP:conf/msr/AhmedDTP25} evaluated how human annotations can be replaced by LLM generated labels.
In their study, they explicitly report the calculated agreement metrics between the models and between humans. They found that three LLMs (GPT-4, Gemini, and Claude) had better agreement on some tasks than human annotators. Meanwhile, \citeauthor{DBLP:conf/icse/XueCBTH24}~\cite{DBLP:conf/icse/XueCBTH24} conducted a controlled experiment to evaluate the impact of ChatGPT on the performance and perceptions of students in an introductory programming course.
They used multiple measures to judge LLM impacts from the human point of view including recording students' screens, evaluating their answers for given tasks, and distributing an opinion survey.

%\citeauthor{hymel2025analysisllmsvshuman}~\cite{hymel2025analysisllmsvshuman} evaluated the capability of ChatGPT-4.0 to generate requirements documents. 
%Specifically, they evaluated two requirements documents based on the same business use case, one document generated with the LLM and one document created by a human expert.
%The documents were then reviewed by experts and judged in terms of alignment with the original business use case (scale of 1-10), requirements completion (Not Complete, Fairly Complete, Fully Complete), and whether they believed it was created by a human or an LLM.
%Finally, they analyzed the influence of the participants' familiarity with AI tools on the study results.
%Participants self-reported their AI tool proficiency (Novice, Intermediate, Advanced, Expert) and usage frequency (Daily, Weekly, Monthly, Never), which were then correlated with their ratings of the requirements documents.

\guidelinesubsubsection{Benefits}

Validating LLMs against human judgments builds confidence in the LLM's accuracy and validity, increasing the trustworthiness of findings, especially when IRA or IRR metrics are reported~\cite{khraisha2024canlargelanguagemodelshumans}. Furthermore, incorporating feedback from human judges may help to improve the LLM or LLM-based tool based on the reported experiences.


%For example, LLM results vary significantly in natural language processing tasks, casting doubt on their reliability and the possibiltiy of replacing human judges~\cite{DBLP:journals/corr/abs-2406-18403}. Moreover, LLMs do not match human annotation in labeling tasks for natural language inference, position detection, semantic change, and hate speech detection~\cite{DBLP:conf/chi/Wang0RMM24}.



\guidelinesubsubsection{Challenges}

Assessing a construct like LLM performance is challenging because ensuring that a construct is defined well and operationalized using an appropriate measurement model requires a deep understanding of (1) the construct, (2) construct validity in general, and (3) instrumentation~\cite{DBLP:journals/tse/SjobergB23,ralph2024teaching}. Comparing an LLM to human judges is typically slower and more expensive than machine-generated measures. More fundamentally, neither human judgment nor machine-generated measures provides an objective ground truth against which LLM accuracy can be firmly determined.   

Human judgments exhibit variability due to differences in experience, expertise, interpretations, and personal biases~\cite{DBLP:journals/pacmhci/McDonaldSF19}. When diverse humans---given clear decision rules corresponding to our construct definitions---rate items reliably, \textit{we simply assume that their reliability implies measurement validity}. It does not. We assume that reliability implies validity because measuring reliability is \textit{much} easier than measuring validity. Indeed, much of the time the best researchers can do is make a conceptual argument for why these judges armed with these decision rules based on this construct \textit{should} produce valid ratings. 


%For example, \citeauthor{hicks_lee_foster-marks_2025} found that racialized software developers \enquote{rated the quality of the output of AI-assisted coding tools significantly lower than other groups}~\cite{hicks_lee_foster-marks_2025}.
% Such biases can be mitigated with a careful selection and combination of multiple judges, but cannot be completely controlled for.

Researchers must weigh the cost and time need for human validation against the expected corresponding improvement in validity over machine-generated measures. 

\guidelinesubsubsection{Study Types}

For \llmusage, researchers \should carefully reflect on the validity of their evaluation criteria.
When an extensively-validated and widely-used benchmark is available (see \benchmarkingtasks), there is little need for human validation.
Where criteria are less clear (e.g., \annotators, \subjects, or \synthesis), or researchers are creating and validating a new benchmark, human validation is important, and this guideline should be followed.
When using \judges, it is common to start with humans who co-create initial rating criteria. 
In developing \newtools, human validation of the LLM output is critical when the output is supposed to match human expectations.

\guidelinesubsubsection{Advice for Reviewers}

\added{Human validation may be the most difficult of our guidelines for reviewers to assess because it will often involve cutting through fallacious arguments. Authors have strong incentives to avoid human validation. However, if LLM output is validated only by comparison with the same or other LLMs, reviewers should demand \textit{quantitative empirical evidence} that such comparison is reliable \textit{and} valid. Similarly, reviewers should demand \textit{quantitative empirical evidence} that any employed benchmarks are reliable \textit{and} valid. Without such evidence, insist on human validation. Showing that several LLMs exhibit high agreement is not sufficient because reliability is necessary but insufficient for validity.}

\added{Next, reviewers will encounter dubious approaches to human validation, such as having a single judge (typically an author) or recruiting dozens of anonymous, untrained, inexperienced gig-workers and aggregating their results on a majority-vote basis. Reviewers should only accept a single judge when the judgments depend only on widely-accepted theories and entails limited value conflict (e.g. tagging method names that contain abbreviations or acronyms). Unfortunately, researchers tend to underestimate the level how theory- and value-laden phenomena are. Any assessment of code quality or tagging of bugs, for example, would be highly dependent on the judge's values and implicit theories~\cite{ralph2024teaching}. }

\added{For multiple judges, reviewers should expect research designs that include techniques known to improve IRA/IRR such as recruiting raters with appropriate education and experience, organizing the rating task into rounds, having raters conduct the initial round of judgments together, having consensus meetings (not voting) after each round in which decision rules are updated. Reviewers should not accept low IRR or IRA (e.g. Krippendorff's $\alpha<0.8$) for research designs absent any of these techniques. Conversely, if authors have embraced best practices for improving IRR/IRA and still obtained middling results (e.g. $0.66<\alpha<0.8$) reviewers should merely insist that low reliability be noted as a limitation.} 

\added{Moving beyond reliability, reviewers should expect authors to \textit{explain} (conceptually, not qualitatively) why human judgments should be valid by comparing construct definitions, directions, decision rules, and judge expertise.}

\added{As with previous guidelines, when important information is missing, reviewers should simply request it be added unless so much is missing that the reviewer cannot assess methodological rigor. For example, failing to provide judge instructions, instruments, decision rules, or construct definitions may prevent assessment of rigor and validity, while missing information about recruitment is less serious. Reviewers should be patient when requesting clearer definitions---the same construct definitions that seem clear to authors often seem opaque to reviewers, and papers should not be rejected over routine clarification requests.}


\subsection{Use an Open LLM as a Baseline}
\label{sec:use-an-open-llm-as-a-baseline}

\vspace{0.5\baselineskip}
\begin{framed}
\underline{tl;dr:} When researchers can control which LLM a study uses, they \should include an open-weight LLM, ideally an open-source LLM, as a baseline and report inter-model agreement. Researchers \should ensure the open-weight baseline is independently reproducible from their \supplementarymaterial.
 
\end{framed}

\guidelinesubsubsection{Rationale}

Reproducibility depends on access to the model under study.
When research relies exclusively on proprietary models, other researchers cannot independently verify or build upon the findings.
Including an open LLM as a baseline ensures that at least part of the study can be fully replicated.

\guidelinesubsubsection{Recommendations}

Empirical studies using LLMs in SE, especially those that target commercial tools or models, \should incorporate an open LLM as a baseline and report established metrics for inter-model agreement (see \benchmarksmetrics).
We acknowledge that including an open LLM baseline might not always be possible, for example, if the study involves human participants, and letting them work on the tasks using two different models might not be feasible.
Using an open model as a baseline is also not necessary if the use of the LLM is tangential to the study goal.

Open models allow other researchers to verify research results and build upon them, even without access to commercial models.
A comparison of commercial and open models also allows researchers to contextualize model performance.
Researchers \should ensure the open-LLM baseline is independently reproducible from their \supplementarymaterial.

Open LLMs are available from hubs such as \href{https://huggingface.co/}{\emph{Hugging Face}}. They can be self-hosted with frameworks such as \href{https://ollama.com/}{\emph{Ollama}} or \href{https://lmstudio.ai/}{\emph{LM Studio}}, accessed through cloud services such as \href{https://together.ai/}{\emph{Together AI}}, AWS, Azure, and Google Cloud, or routed through aggregators such as \href{https://openrouter.ai/}{\emph{OpenRouter}} that expose many providers behind a single API. For agentic setups, open-source tools such as \href{https://www.continue.dev/}{\emph{Continue}}, \href{https://cline.bot/}{\emph{Cline}}, and \href{https://opencode.ai/}{\emph{opencode}}~\cite{opencode} publish their full agent code, system prompts, and tool catalog. By contrast, vendor-hosted services such as \emph{GitHub Copilot} and \emph{Claude Code} expose only parts of their tooling (e.g., editor extensions, hook examples) and keep their deployed agent loops, system prompts, and tool catalogs proprietary.
% GROUNDING opencode: "OpenCode is an open source AI coding agent."

The term ``open'' can have different meanings in the context of LLMs.
\citeauthor{widder2024open} discuss three types of openness (i.e., transparency, reusability, and extensibility) and what openness can and cannot provide~\cite{widder2024open}.
% GROUNDING widder2024open: "We highlight three main affordances of 'open' AI, namely transparency, reusability, and extensibility"
The \emph{Open Source Initiative} (OSI)~\cite{OSIAI2024} defines open-source AI as having access to everything needed to understand, modify, share, retrain, and recreate the model.
% GROUNDING OSIAI2024: "Sufficiently detailed information about the data used to train the system so that a skilled person can build a substantially equivalent system."

\guidelinesubsubsection{Examples}

An increasing number of studies have adopted open LLMs as baseline models.
For example, \citeauthor{DBLP:journals/corr/abs-2406-09834} evaluated seven advanced LLMs, six of which were open-source, testing 145 API mappings drawn from eight popular Python libraries across 28,125 completion prompts aimed at detecting deprecated API usage in code completion~\cite{DBLP:journals/corr/abs-2406-09834}.
% GROUNDING DBLP:journals/corr/abs-2406-09834: "This study involved seven advanced LLMs, 145 API mappings from eight popular Python libraries, and 28,125 completion prompts."
\citeauthor{DBLP:journals/corr/abs-2408-04430} compared four LLMs on a cross-language code clone detection task~\cite{DBLP:journals/corr/abs-2408-04430}.
% GROUNDING DBLP:journals/corr/abs-2408-04430: "We evaluated four LLMs and one embedding model for cross-lingual code clone detection."
Three evaluated models were open-source.
\citeauthor{DBLP:conf/eicc/GoncalvesSC0MPS25} fine-tuned the open LLM \emph{LLaMA} 3.2 on a refined version of the \emph{DiverseVul} dataset to benchmark vulnerability detection performance~\cite{DBLP:conf/eicc/GoncalvesSC0MPS25}.
% GROUNDING DBLP:conf/eicc/GoncalvesSC0MPS25: "we present a refined version of DiverseVul dataset, which is used to fine-tune a large language model, LLaMA 3.2, for vulnerability detection"
\citeauthor{DBLP:journals/corr/abs-2601-11895} evaluated nine code completion models on the \emph{DevBench} benchmark and included three open-weight models (\emph{DeepSeek-V3}, \emph{DeepSeek-V3.1}, and \emph{Ministral-3B}) alongside commercial frontier models, releasing benchmark, evaluation scripts, and per-model raw completions under an MIT license~\cite{DBLP:journals/corr/abs-2601-11895}.
% GROUNDING DBLP:journals/corr/abs-2601-11895: "Claude 4 Sonnet, Claude 3.7 Sonnet, GPT-4.1 mini, GPT-4.1, GPT-4o, DeepSeek-V3, DeepSeek-V3.1, GPT-4.1 nano, Ministral 3B"
\emph{CodeBERT}, a bimodal transformer pre-trained by Microsoft Research, is published with model weights, source code, and data processing scripts on GitHub~\cite{codebert}. It has been used as an open baseline across diverse SE tasks, including exploit code generation~\cite{DBLP:journals/jss/YangZCZHC23}, vulnerability detection~\cite{DBLP:conf/gaiis/XiaSD24}, code clone detection~\cite{DBLP:conf/kbse/SonnekalbGBM22}, and programming assistance for exception handling~\cite{DBLP:conf/icse/CaiYMMN24}.
% GROUNDING codebert: "This repo provides the code for reproducing the experiments in CodeBERT: A Pre-Trained Model for Programming and Natural Languages."
% GROUNDING DBLP:journals/jss/YangZCZHC23: "we propose a novel template-augmented exploit code generation approach ExploitGen based on CodeBERT"
% GROUNDING DBLP:conf/gaiis/XiaSD24: "this study introduces a software defect detection system leveraging CodeBERT and Bi-LSTM"
% GROUNDING DBLP:conf/kbse/SonnekalbGBM22: "Transformer networks such as CodeBERT already achieve outstanding results for code clone detection in benchmark datasets"
% GROUNDING DBLP:conf/icse/CaiYMMN24: "we design Neurex via a multi-tasking model with the fine-tuning of the large language model CodeBERT for the three above exception handling recommending tasks"

\guidelinesubsubsection{Benefits}

A true open LLM baseline improves reproducibility by exposing model architectures, parameter settings, and ideally training data, enabling independent verification of results.
Such baselines also let researchers compare novel methods against a stable reference point, since proprietary models can silently change between tests.
They also allow inspection of training data (when released) and model behavior, helping identify biases and limitations.
Unlike closed-source alternatives, which can be withdrawn or silently updated, open LLMs remain available for future studies.
They typically avoid the per-use API fees that can constrain research groups with limited budgets.

\guidelinesubsubsection{Challenges}

Open-source LLMs face several challenges:
\begin{itemize}
    \item \emph{Definitional inconsistency.} Many models release only trained weights without training data or methodological details (``open weight'' openness)~\cite{Gibney2024}, which is why we reference the OSI definition in our recommendations~\cite{OSIAI2024}.
    % GROUNDING Gibney2024: "Technology giants such as Meta and Microsoft are describing their artificial intelligence (AI) models as 'open source' while failing to disclose important information about the underlying technology"
    % GROUNDING OSIAI2024: "The complete source code used to train and run the system"
    \item \emph{Performance gap.} Open models often lag behind the most advanced proprietary ``frontier'' models in common benchmarks, making it difficult to demonstrate clear improvements when evaluating new methods using open LLMs alone.
    \item \emph{Hardware demands.} Deploying and experimenting with these models typically requires substantial hardware resources, in particular high-performance GPUs that may be beyond reach for many academic groups.
    \item \emph{Operational complexity.} Unlike APIs provided by proprietary vendors (e.g., the OpenAI API), installing, configuring, and fine-tuning open-source models can be technically demanding.
\end{itemize}

\guidelinesubsubsection{Study Types}

This guideline applies primarily to study types in which the researcher controls which LLM is used.
For \benchmarkingtasks, an open LLM \should be one of the models under evaluation, so that the reported scores can be independently re-run. Where this is not feasible, researchers \should justify the omission and, per \design, ensure the evaluation harness can be used with open models.
For controlled experiments under \llmusage, an open LLM \should be one of the models under test; if the experimental design requires capabilities that only a specific commercial model exhibits, researchers \should acknowledge this as a limitation.
When evaluating \newtools, researchers \should use an open LLM as a baseline whenever it is technically feasible; if integration proves too complex, they \should report the initial benchmarking results of open models.
For \annotators and \judges, researchers \should compare annotation or judgment quality from open vs.\ commercial models to assess the extent to which results depend on a specific proprietary model.
For \synthesis, researchers \should compare synthesis results from open and commercial models to evaluate robustness of the findings.
For \llmusage, using an open LLM as a baseline is often not feasible when the study observes participants using specific commercial tools; in such cases, investigators \should explicitly acknowledge its absence and discuss how this limitation might affect their conclusions.
For \subjects, using an open LLM as a baseline may similarly be impractical if the study design requires the capabilities of a specific model; researchers \should acknowledge this limitation when applicable.

\guidelinesubsubsection{Advice for Reviewers}

Reviewers should distinguish between LLM use that is central to the research (e.g., building LLM-driven SE tools, using LLMs for synthesis) and use that is tangential (e.g., generating recruiting materials). An open LLM baseline is expected only when LLM use is central; otherwise, its absence need not be justified.
When an open baseline is expected, reviewers should look for either the use of an open LLM or a convincing argument for why it is impractical. If authors claim openness, some justification for that characterization is appropriate. Where practical, reviewers should examine whether the replication package contains sufficiently detailed instructions.
Performance differences between open and proprietary models are beyond the study authors' control. Reviewers should focus on whether the open baseline serves the study's methodological purpose rather than on its absolute performance.

\guidelinesubsubsection{See Also}
\begin{itemize}[label=$\rightarrow$]
    \item \seemodelversion: Authors must report version, configuration, and parameters for open models, not just commercial ones.
    \item \seedesign: Self-hosting an open model adds hardware, hosting, and infrastructure to document.
    \item \seebenchmarksmetrics: Inter-model agreement metrics make open-vs-commercial comparisons interpretable.
    \item \seelimitationsmitigations: When using an open LLM as a baseline is impractical, authors must report this as a limitation.
\end{itemize}



\subsection{Use Suitable Baselines, Benchmarks, and Metrics}
\label{sec:use-suitable-baselines-benchmarks-and-metrics}

\vspace{0.5\baselineskip}
\begin{framed}
\underline{tl;dr:} \input{_tldr/07_use-baselines-benchmarks-and-metrics-tldr.tex}  
\end{framed}

\guidelinesubsubsection{Recommendations}

Benchmarks, baselines, and metrics play an important role in assessing the effectiveness of LLMs and LLM-based tools. \added{A \textit{benchmark} is a ``standard tool for the competitive evaluation and comparison of competing systems or components according to specific characteristics, such as performance, dependability, or security''~\cite{Kistowski2015Benchmark} }
Benchmarks can be used to assess LLM performance on specific tasks such as code summarization, code generation, \added{and static analysis}.
Meanwhile, ``a \textit{metric} is a method, algorithm, or procedure for assigning one or more numbers to a phenomenon''~\cite{ralph2024teaching}. 
A baseline represents a reference point for the measured LLM.
Since LLMs require substantial hardware resources, baselines allow research to compare LLMs against traditional algorithms with lower computational costs.

When selecting benchmarks and metrics, it is important to fully understand the benchmark tasks, what exactly is being measured, and how it relates to the (often latent) variables researchers actually care about. If one or more metrics or benchmarks are used, researchers:
\begin{itemize}
    \item \must briefly justify (in the \paper) why selected benchmarks and metrics are suitable for the given task or study. 
    % \item \must discuss the reliability and validity (especially construct validity) of selected metrics and benchmarks
    \item \should summarize the structure and tasks of the selected benchmark(s), including the programming language(s) and descriptive statistics such as the number of contained tasks and test cases;
    \item \should discuss the limitations of the selected benchmark(s) (e.g. widely used benchmarks such as \emph{HumanEval} \cite{DBLP:journals/corr/abs-2107-03374} and \emph{MBPP} \cite{DBLP:journals/corr/abs-2108-07732} Python functions assess a very specific part of software development, which is not representative of the full breadth of SE work~\cite{Chandra2025benchmarks}).
    \item \added{\should} include an example of a task and the corresponding test case(s) to illustrate the structure of the benchmark.
\end{itemize}

If multiple benchmarks exist for the same task, researchers \should compare performance between benchmarks.
When selecting only a subset of all available benchmarks, researchers \should use the most specific benchmarks given the context.

Furthermore, researchers \should check whether a less resource-hungry approach (e.g., for static analysis tasks or program repair) can serve as a baseline. If so, the LLM or LLM-based tool should be compared to such baselines using suitable metrics. Even if LLM-based tools outperform baselines, researchers \should discuss whether the resources consumed justify the (often marginal) improvements~\cite{DBLP:journals/cacm/Menzies25}.


To compare traditional and LLM-based approaches or different LLM-based tools, researchers \should report established metrics whenever possible, as this allows secondary research.
They can, of course, report additional metrics that they consider appropriate.

In the following, we briefly discuss two common metrics used for generation tasks: \emph{BLEU-N} and \emph{pass@k}.
\emph{BLEU-N} \cite{DBLP:conf/acl/PapineniRWZ02} is a similarity score based on n-gram precision between two strings, ranging from $0$ to $1$.
Values close to $0$ represent dissimilar content, values closer to $1$ represent similar content, indicating that a model is more capable of generating the expected output.
%The score measures how many n-grams from the generated output of a model appear in the target string.
%In isolation, this approach favors shorter over longer generated sequences because it is more likely to have fewer n-grams in an extensive target sequence than more.
%To mitigate this, the so-called brevity penalty is introduced to disincentivize short sequences when having longer target sequences.
\emph{BLEU-N} has multiple variations.
\emph{CodeBLEU} \cite{DBLP:journals/corr/abs-2009-10297} and \emph{CrystalBLEU} \cite{DBLP:conf/kbse/EghbaliP22} are the most notable ones tailored to code.
They introduce additional heuristics such as AST matching.
As mentioned, researchers \must motivate why they chose a certain metric or variant thereof for their particular study.

The metric \emph{pass@k} reports the likelihood that a model correctly completes a code snippet at least once within \emph{k} attempts.
To our knowledge, the basic concept of \emph{pass@k} was first reported by \citeauthor{DBLP:journals/corr/abs-1906-04908} to evaluate code synthesis~\cite{DBLP:journals/corr/abs-1906-04908}.
They named the concept \emph{success rate at B}, where B denotes the trial ``budget.''
The term \emph{pass@k} was later popularized by \citeauthor{DBLP:journals/corr/abs-2107-03374} as a metric for code generation correctness~\cite{DBLP:journals/corr/abs-2107-03374}.
The exact definition of correctness varies depending on the task.
For code generation, correctness is often defined based on test cases: A passing test implies that the solution is correct.
The resulting pass rates range from $0$ to $1$.
A pass rate of $0$ indicates that the model was unable to generate a single correct solution within $k$ tries; a rate of $1$ indicates that the model successfully generated at least one correct solution in $k$ tries.

The \emph{pass@k} metric is defined as:

\begin{equation}
\text{pass@}k = 1 - \frac{\binom{n-c}{k}}{\binom{n}{k}}
\end{equation}

where $n$ is the total number of generated samples per prompt, $c$ is the number of correct samples among  $n$, and $k$ is the number of samples.

Choosing an appropriate value for \emph{k} depends on the downstream task of the model and how end-users interact with it.
% k= 1 -- pass@1
A high pass rate for \emph{pass@1} is highly desirable in tasks where the system presents only one solution or if a single solution requires high computational effort.
For example, code completion depends on a single prediction, since end users typically see only a single suggestion.
%In automated code refactoring, a large context might lead to high computational costs, making the calculation of a single solution expensive.
%Thus, having a high pass rate at pass@1 is crucial for such tasks.
% k = >=2 -- pass@{>=2}
Pass rates for higher $k$ values (e.g., $2$, $5$, $10$) indicate how well a model can solve a given task in multiple attempts.
%For downstream tasks that allow multiple solutions, e.g. via user interaction, strong performance at $k > 1$ can be justified. 
%For instance, a user selecting the correct test case from multiple suggestions allows for some model errors.
%In code search, finding the most appropriate good snippet for a given context, a reranking algorithm mitigate low pass rates at low k-values.
The \emph{pass@k} metric is frequently referenced in papers on code generation \cite{DBLP:journals/corr/abs-2308-12950, DBLP:journals/corr/abs-2401-14196, DBLP:journals/corr/abs-2409-12186, DBLP:journals/corr/abs-2305-06161}.
However, \emph{pass@k} is not a universal metric suitable for all generation tasks.
It requires a binary notion of correctness, making the metric appropriate for code synthesis evaluated via unit tests, but unsuitable for open-ended generation tasks such as comment generation, where multiple valid outputs exist.

If a study evaluates an LLM-based tool for supporting humans, a relevant metric is the acceptance rate, meaning the ratio of all accepted artifacts (e.g., test cases, code snippets) in relation to all artifacts that were generated and presented to the user.
Another way of evaluating LLM-based tools is calculating inter-model agreement (see also Section \href{/guidelines/#use-an-open-llm-as-a-baseline}{Use an Open LLM as a Baseline}).
This allows researchers to assess how dependent a tool's performance is on specific models and versions.
Metrics used to measure inter-model agreements include general agreement (percentage), \emph{Cohen's kappa}, and \emph{Scott's Pi coefficient}.

Due to LLM non-determinism, researchers \should repeat experiments to  assess (statistically) model or tool performance, e.g., using the arithmetic mean, median, confidence intervals, standard deviations, or more advanced statistical approaches~\cite{DBLP:conf/nips/AgarwalSCCB21, bjarnason2026randomnessagenticevals}.

From a measurement perspective, researchers \should reflect on the theories, values, and measurement models on which the benchmarks and metrics they have selected for their study are based.
For example, the phenomena labeled ``bugs'' in a large open dataset of software bugs is according to a certain theory of what constitutes a bug, and from the values and perspectives of the people who labeled the dataset. Reflecting on the context in which these labels were assigned and discussing whether and how the labels generalize to a new study context is crucial.

\guidelinesubsubsection{Example(s)}

According to \citeauthor{10.1145/3695988}, the main problem types for LLMs are classification, recommendation, and generation, each of which requires different metrics. In their systematic review,
\citeauthor{hu2025assessingadvancingbenchmarksevaluating} \cite{hu2025assessingadvancingbenchmarksevaluating} categorized 191 LLM benchmarks according to the specific SE task they address, making the paper a valuable resource for identifying existing metrics. 
Metrics used to assess generation tasks include \emph{BLEU}, \emph{pass@k}, \emph{Accuracy}, \emph{Accuracy@k}, and \emph{Exact Match}.
The most common metric for recommendation tasks is \emph{Mean Reciprocal Rank}.
For classification tasks, classical machine learning metrics such as \emph{Precision}, \emph{Recall}, \emph{F1-score}, and \emph{Accuracy} are often reported.

%\guidelinesubsubsection{Exact Match}
%\emph{Exact match} is another metric that calculates the percentage of replicas a model can produce.
%If the model is able to produce the exact target sequence, a score of 1 is awarded; otherwise 0.
%Compared to BLEU-N, \emph{exact match} is a stricter measurement.
%Due to its strictness, reported scores are usually low.

%he choice of metrics strongly depends on the specific task.
%Consequently, papers that address different research questions employ different metrics.
%For example, \citeauthor{DBLP:conf/nips/LiuXW023} assesses LLM models for their correctness in code generation using \emph{pass@k}~\cite{DBLP:conf/nips/LiuXW023}.
%However, \emph{pass@k} is not an universal metric suitable for every problem at hand.
%For instance, for creative or open-ended tasks such as comment generation, pass@k is not a suitable metric since not only one single correct result exists.
%Thus, \citeauthor{DBLP:conf/icse/GengWD00JML24} evaluates LLMs for code comment generation using BLEU, ROUGE-L and METEOR as evaluation %metrics~\cite{DBLP:conf/icse/GengWD00JML24}.
%\citeauthor{DBLP:journals/ase/LiuJZNLL25} investigates the performance of LLMs on automated refactorings tasks and proposes a novel metric tailored to the unique characteristics of this problem \cite{DBLP:journals/ase/LiuJZNLL25}.

%Code snippets provided by benchmarks contain untrusted code, making sandboxing recommendable.
%To avoid potential security risks, execute untrusted code in an isolated environment, such as a virtual machine or containerization, instead of your own environment.

Benchmarks used for code generation incldue \emph{HumanEval} (available on \href{https://github.com/openai/human-eval}{GitHub}) \cite{DBLP:journals/corr/abs-2107-03374}, \emph{MBPP} (available on \href{https://huggingface.co/datasets/google-research-datasets/mbpp}{Hugging Face}) \cite{DBLP:journals/corr/abs-2108-07732},
% Both benchmarks consist of code snippets written in Python sourced from publicly available repositories.
% Each snippet consists of four parts: a prompt containing a function definition and a corresponding description of what the function should accomplish, a canonical solution, an entry point for execution, and test cases.
% The input of the LLM is the entire prompt.
% The output of the LLM is then evaluated against the canonical solution using metrics or against a test suite.
\emph{ClassEval} (available on \href{https://github.com/FudanSELab/ClassEval}{GitHub}) \cite{DBLP:journals/corr/abs-2308-01861}, \emph{LiveCodeBench} (available on \href{https://github.com/LiveCodeBench/LiveCodeBench}{GitHub}) \cite{DBLP:journals/corr/abs-2403-07974}, and \emph{SWE-bench} (available on \href{https://github.com/swe-bench/SWE-bench}{GitHub}) \cite{DBLP:conf/iclr/JimenezYWYPPN24}.
An example of a code translation benchmark is \emph{TransCoder}~\cite{DBLP:journals/corr/abs-2006-03511} (available on \href{https://github.com/facebookresearch/CodeGen}{GitHub}). 

% Lukas Boehme: Add examples for other categories: classification and recommendation
Given the generative nature of LLMs, most benchmarks focus on generation tasks, but benchmarks for classification and recommendation tasks also exist.
For classification tasks such as vulnerability detection, \emph{DiverseVul} (available on \href{https://github.com/wagner-group/diversevul}{GitHub})~\cite{DBLP:conf/raid/0001DACW23} provides a dataset of vulnerable and non-vulnerable functions.
They evaluate the binary classification using  \emph{Accuracy}, \emph{Precision}, \emph{Recall}, \emph{F1-score} and \emph{False Positive Rate}.
For recommendation tasks such as code search, \emph{CodeSearchNet} (available on \href{https://github.com/github/CodeSearchNet}{GitHub})~\cite{DBLP:journals/corr/abs-1909-09436} contains code-documentation pairs across multiple programming languages.
The recommendations are evaluated using \emph{Mean Reciprocal Rank}.

\guidelinesubsubsection{Benefits}

Benchmarks and metrics improve reproducibility and comparability across studies, leading to faster improvements in the cumulative body of knowledge. It is simply easier to assess a new system against one or a few benchmarks than hundreds of competing systems, and when most systems are assessed against the same benchmarks and metrics, we can see which approaches work best (at least, on those benchmarks). This comparison enables progress tracking; for example, when researchers iteratively improve a new LLM-based tool and test it against benchmarks after significant changes. For practitioners, leaderboards (published benchmark results of models) support selecting models for downstream tasks. \added{Meanwhile, baselines help us understand just how much LLMs improve performance over non-LLM-enabled alternatives.} 

\guidelinesubsubsection{Challenges}

\added{Computer scientists have a long history of inventing oversimplified, unvalidated metrics (e.g. lines of code, CPU time) and simply claiming that they are reasonable proxies for complicated, value laden, under-theorized, often multidimensional, latent variables (e.g. system size, environmental sustainability). In scientific fields that take measurement more seriously, researchers expect that metrics and instruments are backed up by foundational theories (e.g. we accept telescopes based on our understanding of the physics of light) or empirical evidence (e.g. we accept questionnaire scales for the Big Five personality traits based on extensive quantitative and qualitative investigations of their psychometric properties). The fundamental challenge with AI metrics and benchmarks is that many of them have neither firmly understood theoretical underpinnings nor extensive empirical validation of their construct, measurement, and ecological validity.} 

For example, metrics such as \emph{pass@k} do not reflect qualitative aspects of the generated source code, including its maintainability or readability.
However, these aspects are critical for many downstream tasks.
Moreover, benchmarks are isolated test sets and may not fully represent real-world applications.
For example, benchmarks such as \emph{HumanEval} synthesize code based on written specifications.
However, such explicit descriptions are rare in real-world applications.
Thus, evaluating model performance with benchmarks might not reflect real-world tasks and end-user performance.

Benchmarks and metrics also have generalizability problems. Prominent LLM benchmarks such as \emph{HumanEval} and \emph{MBPP} use Python, so researchers can optimize for Python's idiosyncrasies, creating illusory improvements that do not generalize to other languages. Similarly, the metric \emph{BLEU-N} is a syntactic metric, so code can score highly without being executable. The metric \emph{Exact Match}, meanwhile, does not account for functional equivalence of syntactically different code. Both \emph{BLEU-N} and \emph{Exact Match} are influenced by code formatting, which confounds their intended use. 

Execution-based metrics such as \emph{pass@k} directly evaluate correctness by running test cases, but they require a setup with an execution environment.
When researchers observe unexpected values for certain metrics, the specific results should be investigated in more detail to uncover potential problems.


% Many closed-source models, such as those released by \emph{OpenAI}, achieve exceptional performance on certain tasks but lack transparency and reproducibility~\cite{DBLP:conf/nips/00110ZZDJLHL24, DBLP:journals/corr/abs-2308-01861, DBLP:journals/corr/abs-2406-15877}.
% Benchmark leaderboards, particularly for code generation, are led by closed-source models~\cite{DBLP:journals/corr/abs-2308-01861, DBLP:journals/corr/abs-2406-15877}.
% While researchers \should compare performance against these models, they must consider that providers might discontinue them or apply undisclosed pre- or post-processing beyond the researchers' control (see \openllm).
%The lack of transparency introduces challenges in analyzing mistakes, limiting the understanding of model failures.
%Open-source models offer greater transparency since the entire process is under the researcher's control; however, they require appropriate hardware.

\added{Finally, benchmark data contamination---where the benchmark itself is part of the training data---may lead to artificially high performance if the model remembers the solution from the training data rather than actually solving the task based on the input data~\cite{DBLP:journals/corr/abs-2406-04244}. Therefore, for proprietary LLMs that do not release their training data, researchers \should consider using human validation (Section~\ref{sec:use-human-validation-for-llm-outputs}), curating new data, or refactoring existing data, or~\cite{DBLP:journals/corr/abs-2406-04244}.} 

\added{Creating and curating new uncontaminated benchmark datasets produces unseen benchmark data.
This can be achieved by collecting data after a specified cutoff date, ensuring that models trained before this date could not have been exposed to the benchmark. To ensure transparency, researchers \must disclose the data sources and their collection dates for each benchmark release. However, this temporal approach may be impractical because it requires continuous updates as new models might include the benchmark in future training datasets. Keeping the benchmark private prevents inclusions in training sets, yet it requires trust in the benchmark creator and a system to execute the benchmark without leaking the benchmark.}

\added{Refactoring existing data restructures benchmarks by generating new prompts or augmenting samples across multiple dimensions. Thus, pure memorization by the model is ineffective. However, this only partially addresses contamination since the model may have learned underlying concepts from other versions of the benchmark.}


%\removed{Models may become overfitted to the training data, leading to poor performance on unseen data. When optimizing for a metric, the underlying model might adjust its parameters to improve the metric, thereby failing to generalize. This can result in a model performing exceptionally well on a benchmark but struggling to hold up the performance in unforeseen settings, such as real-world scenarios.}

\guidelinesubsubsection{Study Types}

This guideline \must be followed for all study types that automatically evaluate the performance of LLMs or LLM-based tools.
The design of a benchmark and the selection of appropriate metrics are highly dependent on the specific study type and research goal.
Recommending specific metrics for specific study types is beyond the scope of these guidelines, but \citeauthor{hu2025assessingadvancingbenchmarksevaluating} provide a good overview of existing metrics for evaluating LLMs~\cite{hu2025assessingadvancingbenchmarksevaluating}.
In addition, our Section \humanvalidation provides an overview of the integration of human evaluation in LLM-based studies. 
For \annotators, the research goal might be to assess which model comes close to a ground truth dataset created by human annotators.
Especially for open annotation tasks, selecting suitable metrics to compare LLM-generated and human-generated labels is important.
In general, annotation tasks can vary significantly.
Are multiple labels allowed for the same sequence? Are the available labels predefined, or should the LLM generate a set of labels independently?
Due to this task dependence, researchers \must justify their metric choice, explaining what aspects of the task it captures together with known limitations.
If researchers assess a well-established task such as code generation, they \should report standard metrics such as \emph{pass@k} and compare the performance between models.
If non-standard metrics are used, researchers \must state their reasoning.

\guidelinesubsubsection{Advice for Reviewers}

\added{Reviewers should expect manuscripts to: 
\begin{enumerate}
    \item clearly identify the constructs or variables that the study ought to measure in principle (e.g. LLM performance, quality of generated code), including independent variable(s), dependent variable(s) and (both physical and statistical) control variables;
    \item present their measurement model; that is, identify exactly which metrics, benchmarks, or baselines are used and how they relate to constructs or variables the study aims to measure;
    \item explicitly justify the selection of any metrics, benchmarks, and baselines used
    item contemplate \textit{in detail} the assumptions, reliability, and validity---\textit{especially construct validity} of each benchmark and metric;
    \item clearly articulate any limitations regard construct and measurement validity (see also Section~\ref{sec:report-limitations-and-mitigations}).
\end{enumerate}}

\added{As in previous guidelines, if some specific information about a baseline, etc. is missing, reviewers should simply ask authors to add it. In contrast, reviewers must not tolerate vague descriptions that conflate what researchers and practitioners actually care about (e.g. effectiveness, efficiency, quality, etc.) with specific methods of counting something ostensibly indicative of these broader concerns. Reviewers should be wary of popularity arguments; while the whole point of benchmarks is that standardizing evaluation is beneficial, ubiquity does not imply validity or appropriateness for any given context. Similarly, reviewers must not tolerate ambivalence toward construct validity or obfuscation of validity concerns.  In summary, for this guideline, reviewers should look for manuscripts to convey a solid understanding of construct and measurement validity by explaining and justifying selected measurement models. }


\subsection{Report Limitations and Mitigations}
\label{sec:report-limitations-and-mitigations}

\vspace{0.5\baselineskip}
\begin{framed}
\underline{tl;dr:} \input{_tldr/08_report-limitations-and-mitigations-tldr.tex}   
\end{framed}

\guidelinesubsubsection{Rationale}

When using LLMs for empirical studies in SE, researchers face unique challenges and potential limitations that can influence the validity, reliability, and reproducibility of their findings~\cite{DBLP:conf/icse/SallouDP24}.
% GROUNDING DBLP:conf/icse/SallouDP24: "researchers must still be careful as numerous intricate factors can influence the outcomes of experiments involving LLMs"
Researchers must openly discuss these limitations and explain how their impact was mitigated.
These limitations are relative to current LLM capabilities and tool architectures; speculating about future improvements is beyond the scope of a paper's limitation section.
Nevertheless, risk management and threat mitigation \should be planned during study design, not as an afterthought.

\guidelinesubsubsection{Recommendations}
Researchers \must clearly present the limitations of their work without defensiveness or obfuscation. These limitations may concern external, internal, and construct validity; reliability and reproducibility; and ethical, regulatory, and environmental concerns.

We follow the standard SE convention of organizing limitations by external, internal, and construct validity, plus reliability~\cite{DBLP:journals/ese/RunesonH09,ralph2021empiricalstandardssoftwareengineering}, extended below with ethical, regulatory, and environmental concerns specific to LLM research. For studies that adopt qualitative analytic methods (the use of LLMs in reflexive qualitative analysis is itself contested, see \annotators), researchers \should use the trustworthiness criteria of credibility, transferability, dependability, and confirmability instead~\cite{guba1981criteria}. When deterministic reproducibility is structurally unattainable (e.g., SaaS-based models with opaque versioning), researchers \should adopt the same trustworthiness criteria to substantiate dependability and confirmability of findings.
% GROUNDING DBLP:journals/ese/RunesonH09: "This scheme distinguishes between four aspects of the validity, which can be summarized as follows: Construct validity ... Internal validity ... External validity ... Reliability"
% GROUNDING ralph2021empiricalstandardssoftwareengineering: "General Quality Criteria (e.g. internal validity, construct validity, credibility, resonance)"
% GROUNDING guba1981criteria: "The four terms credibility, transferability, dependability and confirmability were established as the naturalist's equivalents for the more conventional rigor criteria of internal validity, external validity, reliablity and objectivity."

\paragraph{External Validity.}
The primary threats to external validity are:
\begin{itemize}
    \item \emph{Cross-model transfer limitations}: Results obtained with one LLM or family of LLMs may not generalize to others due to differences in training data, architecture, and post-training procedures (e.g., fine-tuning and reinforcement learning from human feedback); see \openllm for using an open LLM as a comparison baseline.
    \item \emph{Configuration sensitivity}: Results may not generalize beyond the specific configuration(s) tested (e.g., decoding parameters, system prompt, or context settings); see \design for an overview.
    \item \emph{Tool-architecture specificity}: Tools built around vendor-specific APIs or features (e.g., function calling, structured output, or context-window size) may not transfer to other models without substantial re-engineering (see \design).
    \item \emph{Limited domain coverage}: Studies often focus on a narrow set of programming languages, task types, or application domains, limiting generalizability to other SE contexts.
    \item \emph{Limited participant and rater diversity}: Study participants (e.g., developers) and human raters validating LLM outputs may not represent the broader population in terms of expertise, geographic location, or cultural background (see \humanvalidation for guidance on rater diversity in value-laden constructs).
    \item \emph{Research-to-practice gap}: Developers using the same tools outside researcher-supervised conditions may obtain results that differ from those reported in the study.
    \item \emph{Cross-time instability}: Performance of proprietary models can change over time, leading to non-generalizable study outcomes~\cite{DBLP:journals/corr/abs-2307-09009, doi:10.1148/radiol.232411} (see \modelversion for version and fingerprint reporting that enables tracking such drift).
    % GROUNDING DBLP:journals/corr/abs-2307-09009: "the behavior of the "same" LLM service can change substantially in a relatively short amount of time"
    % GROUNDING doi:10.1148/radiol.232411: "Our results suggest performance drift between the March and June snapshots of GPT-3.5 and GPT-4."
\end{itemize}

Generalizability is particularly critical for proprietary and non-deterministic systems whose behavior is subject to drift (i.e., silent changes in model output over time). Researchers \must discuss the limitations and mitigations of external validity.
Mitigations include \emph{triangulation} across multiple models (e.g., proprietary and open), independent datasets, and complementary metrics; \emph{sensitivity analysis} that varies LLM configurations, prompts, architecture decisions, datasets, and where applicable participant backgrounds; and performing \emph{longitudinal re-runs} (see \emph{Reliability \& Reproducibility} below), which also helps detect cross-time instability.

\paragraph{Internal Validity.}
The primary threats to internal validity are:
\begin{itemize}
	\item \emph{Data leakage and contamination}: Inter-dataset duplication can produce training-evaluation overlap, yielding overly optimistic results (see \benchmarksmetrics).
	\item \emph{Evaluation data entering model-improvement pipelines}: Evaluation samples can unintentionally feed retraining or fine-tuning, especially in longitudinal studies involving LLMs.
	\item \emph{Incomplete architecture, prompt, or pipeline reporting}: Undisclosed components introduce hidden confounders (see \design).
\end{itemize}

As transparency on training data is limited for LLMs, researchers \must discuss potential data leakage effects and their impact on results. Concrete mitigations for contamination (e.g., post-cutoff benchmark construction, held-out subsets, canary strings) are discussed in \benchmarksmetrics.

\paragraph{Construct Validity.}
The primary threats to construct validity are:
\begin{itemize}
	\item \emph{Metric-construct mismatch}: Traditional metrics such as BLEU or ROUGE may miss SE-specific aspects such as functional correctness or behavioral equivalence (see \benchmarksmetrics).
	\item \emph{Construct under-specification}: If a construct lacks an operational definition, neither automated metrics nor human raters can apply it consistently.
	\item \emph{Reliability without validity}: High inter-rater or inter-model agreement does not imply that the measurement captures the intended construct; a reliable LLM can be reliably inaccurate (see \humanvalidation).
	\item \emph{Benchmark scope limitations}: Benchmarks commonly ignore runtime behaviors, security implications, readability, testability, and maintainability, yielding results that may not transfer to realistic development settings.
	\item \emph{Capability confounding}: Benchmark performance can blend the target capability with unrelated capabilities such as output format compliance, instruction following, or prompt format sensitivity, inflating apparent scores (see \benchmarksmetrics).
	\item \emph{Over-reliance on benchmark-specific metrics}: Optimizing for a single benchmark may produce dataset-specific shortcuts that pass the benchmark without exhibiting the capability it was designed to test, overstating real-world utility.
	\item \emph{Prompt sensitivity}: Small changes in instructions, formatting, or in-context examples can substantially shift what the LLM appears to measure, making the operationalized construct unstable across prompt variants (see \design).
	\item \emph{Judge biases}: LLM and human judges exhibit systematic biases such as position bias, verbosity bias, and preferences for richly formatted or citation-rich outputs regardless of correctness (see \humanvalidation).
\end{itemize}

If constructs are based on subjective interpretations, purely automated metrics are insufficient. Researchers \must discuss how they ensured quality of subjective results, similarly to qualitative research. The primary mitigation is \emph{human validation} of subjective constructs following quality criteria known from qualitative research (see \humanvalidation).

\paragraph{Reliability \& Reproducibility.}
The primary threats to reliability and reproducibility are:
\begin{itemize}
    \item \emph{Non-deterministic outputs}: Identical prompts and configurations can yield different outputs across runs due to factors such as floating-point arithmetic, batching, and stochastic decoding strategies.
    \item \emph{Infrastructure dependence}: Results may vary depending on the hardware, software stack, and hosting environment used; vendor-imposed quotas, throttling, or pricing changes can further prevent re-running experiments at the original scale, making exact replication challenging across different infrastructure setups.
    \item \emph{Resource inequality}: LLM research is resource-intensive and remains predominantly in the domain of private companies or well-funded research institutions~\cite{schwartz2020green, ahmed2023industry}, excluding researchers from under-resourced institutions.
    % GROUNDING schwartz2020green: "the financial cost of the computations can make it difficult for academics, students, and researchers, in particular those from emerging economies, to engage in deep learning research"
    % GROUNDING ahmed2023industry: "industry increasingly dominates the three key ingredients of modern AI research: computing power, large datasets, and highly skilled researchers"
\end{itemize}

Researchers \must discuss the measures taken to increase reliability and reproducibility.
However, non-deterministic reproducibility is not inherently disqualifying. The trustworthiness criteria introduced above apply particularly to SaaS-based LLM research, where providers frequently deprecate model versions without guaranteeing stable behavior.
Mitigations include providing \emph{replication packages} that cover prompt and architecture specifications, model outputs, and representative examples for partial replicability (ideally accompanied by an implementation using an open model for long-term stability), and performing \emph{longitudinal re-runs} with statistical analyses. When deterministic reproduction is structurally impossible, researchers \should consider \emph{methodological trustworthiness measures} such as triangulation, reflexivity, audit trails, and peer debriefing as complementary measures.

\paragraph{Ethical \& Regulatory Boundaries.}
The primary concerns for ethical and regulatory matters are:
\begin{itemize}
    \item \emph{Use of sensitive or proprietary data}: Studies involving proprietary code, confidential business data, or personally identifiable information may face restrictions on data sharing that limit reproducibility.
    \item \emph{Jurisdictional obligations}: Data protection regulations such as GDPR or CCPA and institutional policies may impose constraints on data collection, processing, and sharing in LLM-based studies.
    \item \emph{Implicit model bias}: Especially for qualitative research, LLMs might \enq{reinforce dominant paradigms and biases} and \enq{identify, replicate and reinforce dominant language and patterns}~\cite{jowsey2025reject} (see \humanvalidation).
    % GROUNDING jowsey2025reject: "Failure to recognize these limitations of GenAI risks analyses that reinforce dominant paradigms and biases."
\end{itemize}

Studies involving sensitive data \must discuss data governance mechanisms tailored toward LLM environments, compliant with applicable jurisdictional obligations. If applicable, researchers \should discuss how model biases potentially impact the study outcomes and how those biases were evaluated.

\paragraph{Environmental \& Sustainability Constraints.}
The primary environmental and sustainability concerns are:
\begin{itemize}
    \item \emph{Energy consumption}: With growing model size, the environmental impact of experiments with LLMs increases, and the substantial energy costs of LLM experiments warrant consideration in study design~\cite{DBLP:conf/acl/StrubellGM19}.
    % GROUNDING DBLP:conf/acl/StrubellGM19: "These accuracy improvements depend on the availability of exceptionally large computational resources that necessitate similarly substantial energy consumption."
    \item \emph{Trade-off between repetition and sustainability}: Repeating experiments increases reliability but also energy consumption, requiring trade-offs during study design.
\end{itemize}

Researchers \should justify the LLM's resource consumption against the benefits over traditional approaches.
Mitigations include \emph{energy preservation} and \emph{cost accounting}. \emph{Energy preservation} involves selecting smaller or newer, less resource-intensive models and applying techniques such as input/output token reduction, model pruning, quantization, or knowledge distillation~\cite{mitu2024hidden} where feasible. Carbon footprint estimation is desirable, but still difficult. \emph{Cost accounting} tracks resource consumption by reporting tokens, service costs, or hardware specifications.
% GROUNDING mitu2024hidden: "Techniques such as model pruning, quantisation, and knowledge distillation can reduce the computational resources needed for training and inference, thereby decreasing the overall carbon footprint"

\guidelinesubsubsection{Examples}

\citeauthor{DBLP:conf/icse/SallouDP24} catalog three categories of LLM-specific threats to validity (i.e., closed-source models, implicit data leakage, and reproducibility) and pair each with concrete mitigation strategies (e.g., versioned model archives, metamorphic test data, multiple replication runs with variability metrics, and detailed execution metadata)~\cite{DBLP:conf/icse/SallouDP24}.
% GROUNDING DBLP:conf/icse/SallouDP24: "potential threats to the validity of LLM-based research including issues such as closed-source models, possible data leakage between LLM training data and research evaluation, and the reproducibility of LLM-based findings"
\citeauthor{DBLP:conf/icse/Du0WWL0FS0L24} pair each threat in their \emph{ClassEval} evaluation with a concrete mitigation: manually constructing the benchmark with multiple annotators to limit data leakage, piloting prompts on held-out tasks to control for prompt sensitivity, and reporting greedy-decoding results to control for non-determinism~\cite{DBLP:conf/icse/Du0WWL0FS0L24}.
% GROUNDING DBLP:conf/icse/Du0WWL0FS0L24: "One potential threat is the data leakage between our benchmark and model training data, thus we manually construct the benchmark ClassEval."

\guidelinesubsubsection{Benefits}

Transparent reporting of limitations and mitigations helps readers calibrate confidence in the findings, makes explicit which threats were addressed and which remain open, and documents design decisions that other authors can borrow or refine. It also keeps a paper's claims proportionate to its evidence.

\guidelinesubsubsection{Challenges}

Identifying limitations one is not already aware of is the hardest part of writing a threats section, particularly for methodological threats outside the team's primary expertise. Publication and reviewing norms can pressure authors to downplay weaknesses, while page limits make exhaustive treatment impractical. Threat lists that recite generic LLM-research issues (e.g., model bias, non-determinism, or contamination) without showing how each one applies to specific design choices in this study leave reviewers unable to tell which risks actually applied.

The threats to validity framework itself is contested within SE. \citeauthor{DBLP:journals/infsof/VerdecchiaELRS23} argue that threats sections too often read as \enq{laundry-lists}~\cite{DBLP:journals/infsof/VerdecchiaELRS23}, \citeauthor{DBLP:conf/esem/LagoRSV24} corroborated this empirically across a decade of ICSE Distinguished Paper Award winners~\cite{DBLP:conf/esem/LagoRSV24}, and \citeauthor{DBLP:journals/tosem/RobillardAEGLNNSSS24} argue for refocusing the discussion on study design trade-offs rather than the standard validity categories~\cite{DBLP:journals/tosem/RobillardAEGLNNSSS24}.
% GROUNDING DBLP:journals/infsof/VerdecchiaELRS23: "TTV sections seem to be mostly formulated as “laundry-lists” of potential threats, lacking a thorough contextualization in the specifics of the study at hand"
% GROUNDING DBLP:conf/esem/LagoRSV24: "the selected papers, despite being selected as best papers in the ICSE conference, reported dominantly shallow or none TTV analysis"
% GROUNDING DBLP:journals/tosem/RobillardAEGLNNSSS24: "First, threats to validity sections are often boilerplate text produced by rote."

\guidelinesubsubsection{Study Types}

Researchers \must follow this guideline for all study types. Transparently reporting limitations and mitigations is a universal requirement, but specific concerns vary by study type.
For \annotators, researchers \must discuss potential biases in label assignment, label reliability limitations, and sensitivity of annotations to prompt wording and model choice.
For \judges, researchers \must address measurement validity concerns, known biases such as position bias or verbosity bias, and the extent to which LLM judgments align with human expert assessments.
For \synthesis, researchers \must discuss the risk of contextual misinterpretation, potential loss of nuance in summarized or aggregated outputs, and reflexivity limitations inherent in using an LLM for qualitative interpretation.
For \subjects, researchers \must discuss the fundamental inability of LLMs to truly simulate human behavior, the risk of stereotype amplification, and the limited ecological validity of simulated responses.
For \llmusage, researchers \must discuss generalizability constraints across different tools and user populations, and acknowledge how observed usage patterns may not transfer to other contexts.
For \newtools, researchers \must discuss replicability constraints arising from dependencies on commercial models, the impact of model updates on tool behavior, and limitations of the evaluation setup.
For \benchmarkingtasks, researchers \must discuss potential data contamination, benchmark scope limitations, capability confounding, and the extent to which benchmark performance generalizes to real-world tasks.

\guidelinesubsubsection{Advice for Reviewers}

Reviewers should verify that the limitation section is comprehensive and appropriate for the specific study type, checking that:
(1) limitations address the specific threats relevant to the study type (e.g., label reliability for annotation studies, simulation fidelity for studies using LLMs as subjects);
(2) mitigations are concrete and correspond to identified limitations rather than being generic statements;
(3) the impact of LLM non-determinism on findings is discussed;
(4) generalizability constraints, across models, configurations, time periods, and populations, are acknowledged.
When important limitations are missing, reviewers should request they be added. The absence of a limitation section, or one that is formulaic or insufficiently specific, is a more serious concern than any individual missing limitation and may warrant a major revision.

\guidelinesubsubsection{See Also}
\begin{itemize}[label=$\rightarrow$]
    \item \seemodelversion: Authors can re-run with different models or configurations to check whether the results depend on those specific choices.
    \item \seedesign: Triangulation across architectures and prompts is one mitigation strategy.
    \item \seetraces: Stored session traces serve as a baseline against which authors can monitor LLM behavior drift over time.
    \item \seebenchmarksmetrics: Benchmark and metric choices are one source of construct-validity threats authors must discuss.
    \item \seeopenllm: An open LLM baseline mitigates cross-model transfer concerns by providing an independently reproducible comparison.
    \item \seehumanvalidation: When automated metrics cannot validly measure a construct, human validation is an alternative.
\end{itemize}


\section{Conclusion}

This paper organizes the landscape of empirical studies involving LLMs in software engineering along seven study types and distills eight guidelines aimed at improving their validity, transparency, and reproducibility.
Our guidelines recommend researchers to 
\begin{enumerate*}[label=(\arabic*)]
\item explicitly declare when and how LLMs are used;
\item report model versions, configuration, and fine-tuning details;
\item describe the complete tool architecture beyond the model;
\item release prompts, their development, and, where possible, interaction logs;
\item validate LLM outputs with humans;
\item include an open LLM as a baseline;
\item select appropriate baselines, benchmarks, and metrics; and
\item report study limitations and mitigations, including costs, potential biases, and environmental impact.
\end{enumerate*}
Adopting these practices will not eliminate LLMs' non-determinism, but it will make claims auditable and results easier to replicate and compare across studies. 
By consistently applying our recommendations, the community can produce more reliable evidence on what LLM can and cannot deliver for software engineering research and practice.
Our guidelines complement existing empirical software engineering guidance while addressing LLM-specific challenges.
Because models and tools evolve, we maintain the guidelines as a living resource and invite the community to refine and extend them: \href{https://llm-guidelines.org/}{llm-guidelines.org}.


\section*{Acknowledgments}

Our study types and guidelines are based on discussion sessions with researchers at the 2024 International Software Engineering Research Network (ISERN) meeting and at the 2nd Copenhagen Symposium on Human-Centered Software Engineering AI (supported by the Carlsberg Foundation with grant CF24-0693 and the Alfred P. Sloan Foundation with grant G-2024-22586 to Daniel Russo).
ChatGPT with GPT-5 was used to generate initial versions of the abstract and conclusion sections of this paper.

%\balance
%\bibliographystyle{IEEEtran}
\bibliographystyle{ACM-Reference-Format}
\bibliography{literature}
%\printbibliography

\end{document}
